\documentclass[../main.tex]{subfiles}
\begin{document}

\chapter{GPU Architecture}
\lessoninfo{
  video={M4, lessons 1--4 (Multi-threaded Architecture; Bank Conflicts; GPU Architecture Pipeline; Global Memory Coalescing)},
  sources={PTX ISA \cite{nvidia2026ptx} (special registers); Harris, \emph{Using Shared Memory in CUDA C/C++} \cite{harris2013sharedmem}; NVIDIA A100 \cite{nvidia2020ampere} and H100 \cite{nvidia2022hopper} whitepapers (HBM figures)},
  textbook={A18 \cite{aamodt2018gpu} ch.~3 (assigned), ch.~4; K12 \cite{kim2012gpgpu} §1.4.1 GPU Pipeline, §1.4.3 GPU Memory Systems; CPG \cite{nvidia2024cudaguide} §4.2 Hardware Multithreading, §5.3.2 Device Memory Accesses; BPG \cite{nvidia2026cudabestpractices} §10.2.1 Coalesced Access to Global Memory, §10.2.3 Shared Memory; PMPP \cite{hwu2022pmpp} ch.~4--6},
  keywords={latency hiding, fine-grained multithreading, context switch, program counter per warp, register file, port, bank, bank conflict, operand collector, scoreboard, active mask, special register, memory coalescing, memory transaction, SRAM, DRAM, HBM},
}

Lesson~1 named the parts of a streaming multiprocessor and Lesson~3
showed how a kernel's threads are packed into blocks and warps and how
many of them an SM can hold.  This lesson opens the SM and follows a
warp instruction through it.  It begins with the reason the SM is
multithreaded at all: a load that misses in the cache would otherwise
stall the pipeline, and a GPU hides that latency by switching among
warps whose program counters and registers are all kept on chip.  The
register file that makes this possible is so wide that it must be split
into banks, which introduces bank conflicts, operand collectors and the
scoreboard that selects a ready warp.  A complete pipeline walk-through
on a small example then shows fetch, decode, register read, execution
and write-back for a warp, including the mask bits that switch lanes
off.  The lesson closes with the memory side: why a warp can generate
32 memory requests from one instruction, what SRAM, DRAM and HBM are,
and how coalescing turns those requests into a few wide transactions.
Lesson~5 turns these mechanisms into optimisation rules, and Lessons~6
and~7 return to the register file and the warp scheduler.

\section{Multi-threaded architecture}
\label{sec:l04-multithreading}

A GPU consists of many SMs, each of which runs its warps by
multithreading and holds a block of shared memory and hardware caches
(the SM diagram of Lesson~1).\source{M4, multithreading; K12 §1.2.1;
A18 ch.~3; CPG §4.2 Hardware Multithreading.}  This section explains
what the multithreading buys, what state it requires, and how that
state differs from a CPU's.

\subsection{Three ways to survive a cache miss}

Consider the four-instruction program of \cref{fig:l04-latency-hiding}
on the five-stage pipeline of Lesson~1 (front end, decode and rename,
schedule, execute, write-back).  The first instruction, \texttt{I1},
loads a word from memory; \texttt{I4} needs the loaded value, while
\texttt{I2} and \texttt{I3} are independent of it.  Suppose the load
suffers a \term{cache miss}[キャッシュミス]: the word is not in the
cache and must come from DRAM, so the execute stage of \texttt{I1}
cannot finish for many cycles.\margin{Latency scale: global-memory
accesses cost \enquote{100s of cycles} on the Tesla and Fermi GPUs of
K12 Table~1.4 \cite{kim2012gpgpu}; an on-chip shared-memory access there
costs 4--32 cycles.}  The pipeline \emph{stalls}: cycles pass in which
no instruction completes.  Three designs react differently.
\begin{description}
  \item[In-order execution] (\cref{fig:l04-latency-hiding}a) issues
    instructions in program order, so \texttt{I2}, \texttt{I3} and
    \texttt{I4} all wait behind the miss, although only \texttt{I4}
    depends on it.
  \item[Out-of-order execution] (\cref{fig:l04-latency-hiding}b) lets
    \texttt{I2} and \texttt{I3} complete during the miss; only
    \texttt{I4} waits.  The gain is limited by the independent
    instructions the hardware can find in one thread.
  \item[Multithreading] (\cref{fig:l04-latency-hiding}c) issues
    \texttt{I1} of four different threads in four consecutive cycles.
    Instructions from different threads are independent by
    construction, so no search for independence is needed, and the four
    misses are outstanding at the same time: four memory requests
    overlap instead of being served one after another.
\end{description}

\begin{figure}[htbp]
  \centering
  % Latency hiding: in-order, out-of-order and multithreaded execution of
% one four-instruction program whose first instruction misses in the
% cache (Lesson 4).  (a) in-order: every later instruction waits behind
% the miss; (b) out-of-order: only the dependent instruction waits;
% (c) multithreading: four threads each issue their load and the four
% misses overlap.  Stages F D S E W; grey = waiting for the miss.
% Redraws the outline figures m4l1_multithreading, m4l1_multithreading_2
% and m4l1_multithreading_3.
% Usage: % Latency hiding: in-order, out-of-order and multithreaded execution of
% one four-instruction program whose first instruction misses in the
% cache (Lesson 4).  (a) in-order: every later instruction waits behind
% the miss; (b) out-of-order: only the dependent instruction waits;
% (c) multithreading: four threads each issue their load and the four
% misses overlap.  Stages F D S E W; grey = waiting for the miss.
% Redraws the outline figures m4l1_multithreading, m4l1_multithreading_2
% and m4l1_multithreading_3.
% Usage: % Latency hiding: in-order, out-of-order and multithreaded execution of
% one four-instruction program whose first instruction misses in the
% cache (Lesson 4).  (a) in-order: every later instruction waits behind
% the miss; (b) out-of-order: only the dependent instruction waits;
% (c) multithreading: four threads each issue their load and the four
% misses overlap.  Stages F D S E W; grey = waiting for the miss.
% Redraws the outline figures m4l1_multithreading, m4l1_multithreading_2
% and m4l1_multithreading_3.
% Usage: \input{figures/l04-latency-hiding}
% Width: body (120 mm); the drawing is about 118 mm wide.
\begin{tikzpicture}[x=2.9mm, y=2.9mm, font=\scriptsize\sffamily]
  % \row{x}{y}{fill}{label}{stages}: one instruction; "-" is a stall cell
  \newcommand{\row}[5]{%
    \node[anchor=east, inner sep=1pt] at (#1-0.3,#2-0.5) {#4};
    \foreach \s [count=\i from 0] in {#5} {
      \ifx\s-
        \draw[fill=ln-rule!40, draw=ln-rule] (#1+\i,#2) rectangle ++(1,-1);
      \else
        \draw[fill=#3] (#1+\i,#2) rectangle ++(1,-1) node[midway] {\s};
      \fi
    }}
  % (a) in-order
  \node[anchor=south west, font=\scriptsize\sffamily\bfseries] at (1,0.3)
    {\iflangja{(a) インオーダ}{(a) in-order}};
  \row{1}{0}{ln-box}{I1}{F,D,S,E,-,-,-,-,-,-,W}
  \row{1}{-1}{ln-box}{I2}{,F,D,S,-,-,-,-,-,-,-,E,W}
  \row{1}{-2}{ln-box}{I3}{,,F,D,-,-,-,-,-,-,-,S,E,W}
  \row{1}{-3}{ln-box}{I4}{,,,F,-,-,-,-,-,-,-,D,S,E,W}
  \draw[->, ln-muted] (1,-4.6) -- (16,-4.6)
    node[below left, inner sep=1pt] {\iflangja{サイクル}{cycle}};
  % (b) out-of-order
  \node[anchor=south west, font=\scriptsize\sffamily\bfseries] at (21,0.3)
    {\iflangja{(b) アウトオブオーダ}{(b) out-of-order}};
  \row{21}{0}{ln-box}{I1}{F,D,S,E,-,-,-,-,-,-,W}
  \row{21}{-1}{ln-box}{I2}{,F,D,S,E,W}
  \row{21}{-2}{ln-box}{I3}{,,F,D,S,E,W}
  \row{21}{-3}{ln-box}{I4}{,,,F,D,S,-,-,-,-,-,E,W}
  \draw[->, ln-muted] (21,-4.6) -- (36,-4.6)
    node[below left, inner sep=1pt] {\iflangja{サイクル}{cycle}};
  % (c) multithreading
  \node[anchor=south west, font=\scriptsize\sffamily\bfseries] at (1,-7.7)
    {\iflangja{(c) マルチスレッディング（各スレッドの I1）}{(c) multithreading (I1 of each thread)}};
  \row{1}{-8}{ln-accent!25}{T1}{F,D,S,E,-,-,-,-,-,-,W}
  \row{1}{-9}{ln-tip!25}{T2}{,F,D,S,E,-,-,-,-,-,-,W}
  \row{1}{-10}{ln-caution!20}{T3}{,,F,D,S,E,-,-,-,-,-,-,W}
  \row{1}{-11}{ln-source!20}{T4}{,,,F,D,S,E,-,-,-,-,-,-,W}
  \draw[->, ln-muted] (1,-12.6) -- (16,-12.6)
    node[below left, inner sep=1pt] {\iflangja{サイクル}{cycle}};
  % program and legend
  \node[anchor=north west, align=left, font=\scriptsize\ttfamily] at (21,-7.7)
    {I1: R1 = LD [0x400]\\I2: R2 = R5 + R6\\I3: R3 = R2 + R6\\I4: R5 = R1 + R6};
  \draw[fill=ln-rule!40, draw=ln-rule] (21,-11.8) rectangle ++(1,-1);
  \node[anchor=west, text=ln-muted] at (22.3,-12.3)
    {\iflangja{ミス待ち（ストール）}{waiting for the miss (stall)}};
  \node[anchor=west, text=ln-muted] at (22.3,-13.5)
    {\iflangja{I4 だけが I1 の R1 に依存する}{only I4 depends on R1 from I1}};
\end{tikzpicture}

% Width: body (120 mm); the drawing is about 118 mm wide.
\begin{tikzpicture}[x=2.9mm, y=2.9mm, font=\scriptsize\sffamily]
  % \row{x}{y}{fill}{label}{stages}: one instruction; "-" is a stall cell
  \newcommand{\row}[5]{%
    \node[anchor=east, inner sep=1pt] at (#1-0.3,#2-0.5) {#4};
    \foreach \s [count=\i from 0] in {#5} {
      \ifx\s-
        \draw[fill=ln-rule!40, draw=ln-rule] (#1+\i,#2) rectangle ++(1,-1);
      \else
        \draw[fill=#3] (#1+\i,#2) rectangle ++(1,-1) node[midway] {\s};
      \fi
    }}
  % (a) in-order
  \node[anchor=south west, font=\scriptsize\sffamily\bfseries] at (1,0.3)
    {\iflangja{(a) インオーダ}{(a) in-order}};
  \row{1}{0}{ln-box}{I1}{F,D,S,E,-,-,-,-,-,-,W}
  \row{1}{-1}{ln-box}{I2}{,F,D,S,-,-,-,-,-,-,-,E,W}
  \row{1}{-2}{ln-box}{I3}{,,F,D,-,-,-,-,-,-,-,S,E,W}
  \row{1}{-3}{ln-box}{I4}{,,,F,-,-,-,-,-,-,-,D,S,E,W}
  \draw[->, ln-muted] (1,-4.6) -- (16,-4.6)
    node[below left, inner sep=1pt] {\iflangja{サイクル}{cycle}};
  % (b) out-of-order
  \node[anchor=south west, font=\scriptsize\sffamily\bfseries] at (21,0.3)
    {\iflangja{(b) アウトオブオーダ}{(b) out-of-order}};
  \row{21}{0}{ln-box}{I1}{F,D,S,E,-,-,-,-,-,-,W}
  \row{21}{-1}{ln-box}{I2}{,F,D,S,E,W}
  \row{21}{-2}{ln-box}{I3}{,,F,D,S,E,W}
  \row{21}{-3}{ln-box}{I4}{,,,F,D,S,-,-,-,-,-,E,W}
  \draw[->, ln-muted] (21,-4.6) -- (36,-4.6)
    node[below left, inner sep=1pt] {\iflangja{サイクル}{cycle}};
  % (c) multithreading
  \node[anchor=south west, font=\scriptsize\sffamily\bfseries] at (1,-7.7)
    {\iflangja{(c) マルチスレッディング（各スレッドの I1）}{(c) multithreading (I1 of each thread)}};
  \row{1}{-8}{ln-accent!25}{T1}{F,D,S,E,-,-,-,-,-,-,W}
  \row{1}{-9}{ln-tip!25}{T2}{,F,D,S,E,-,-,-,-,-,-,W}
  \row{1}{-10}{ln-caution!20}{T3}{,,F,D,S,E,-,-,-,-,-,-,W}
  \row{1}{-11}{ln-source!20}{T4}{,,,F,D,S,E,-,-,-,-,-,-,W}
  \draw[->, ln-muted] (1,-12.6) -- (16,-12.6)
    node[below left, inner sep=1pt] {\iflangja{サイクル}{cycle}};
  % program and legend
  \node[anchor=north west, align=left, font=\scriptsize\ttfamily] at (21,-7.7)
    {I1: R1 = LD [0x400]\\I2: R2 = R5 + R6\\I3: R3 = R2 + R6\\I4: R5 = R1 + R6};
  \draw[fill=ln-rule!40, draw=ln-rule] (21,-11.8) rectangle ++(1,-1);
  \node[anchor=west, text=ln-muted] at (22.3,-12.3)
    {\iflangja{ミス待ち（ストール）}{waiting for the miss (stall)}};
  \node[anchor=west, text=ln-muted] at (22.3,-13.5)
    {\iflangja{I4 だけが I1 の R1 に依存する}{only I4 depends on R1 from I1}};
\end{tikzpicture}

% Width: body (120 mm); the drawing is about 118 mm wide.
\begin{tikzpicture}[x=2.9mm, y=2.9mm, font=\scriptsize\sffamily]
  % \row{x}{y}{fill}{label}{stages}: one instruction; "-" is a stall cell
  \newcommand{\row}[5]{%
    \node[anchor=east, inner sep=1pt] at (#1-0.3,#2-0.5) {#4};
    \foreach \s [count=\i from 0] in {#5} {
      \ifx\s-
        \draw[fill=ln-rule!40, draw=ln-rule] (#1+\i,#2) rectangle ++(1,-1);
      \else
        \draw[fill=#3] (#1+\i,#2) rectangle ++(1,-1) node[midway] {\s};
      \fi
    }}
  % (a) in-order
  \node[anchor=south west, font=\scriptsize\sffamily\bfseries] at (1,0.3)
    {\iflangja{(a) インオーダ}{(a) in-order}};
  \row{1}{0}{ln-box}{I1}{F,D,S,E,-,-,-,-,-,-,W}
  \row{1}{-1}{ln-box}{I2}{,F,D,S,-,-,-,-,-,-,-,E,W}
  \row{1}{-2}{ln-box}{I3}{,,F,D,-,-,-,-,-,-,-,S,E,W}
  \row{1}{-3}{ln-box}{I4}{,,,F,-,-,-,-,-,-,-,D,S,E,W}
  \draw[->, ln-muted] (1,-4.6) -- (16,-4.6)
    node[below left, inner sep=1pt] {\iflangja{サイクル}{cycle}};
  % (b) out-of-order
  \node[anchor=south west, font=\scriptsize\sffamily\bfseries] at (21,0.3)
    {\iflangja{(b) アウトオブオーダ}{(b) out-of-order}};
  \row{21}{0}{ln-box}{I1}{F,D,S,E,-,-,-,-,-,-,W}
  \row{21}{-1}{ln-box}{I2}{,F,D,S,E,W}
  \row{21}{-2}{ln-box}{I3}{,,F,D,S,E,W}
  \row{21}{-3}{ln-box}{I4}{,,,F,D,S,-,-,-,-,-,E,W}
  \draw[->, ln-muted] (21,-4.6) -- (36,-4.6)
    node[below left, inner sep=1pt] {\iflangja{サイクル}{cycle}};
  % (c) multithreading
  \node[anchor=south west, font=\scriptsize\sffamily\bfseries] at (1,-7.7)
    {\iflangja{(c) マルチスレッディング（各スレッドの I1）}{(c) multithreading (I1 of each thread)}};
  \row{1}{-8}{ln-accent!25}{T1}{F,D,S,E,-,-,-,-,-,-,W}
  \row{1}{-9}{ln-tip!25}{T2}{,F,D,S,E,-,-,-,-,-,-,W}
  \row{1}{-10}{ln-caution!20}{T3}{,,F,D,S,E,-,-,-,-,-,-,W}
  \row{1}{-11}{ln-source!20}{T4}{,,,F,D,S,E,-,-,-,-,-,-,W}
  \draw[->, ln-muted] (1,-12.6) -- (16,-12.6)
    node[below left, inner sep=1pt] {\iflangja{サイクル}{cycle}};
  % program and legend
  \node[anchor=north west, align=left, font=\scriptsize\ttfamily] at (21,-7.7)
    {I1: R1 = LD [0x400]\\I2: R2 = R5 + R6\\I3: R3 = R2 + R6\\I4: R5 = R1 + R6};
  \draw[fill=ln-rule!40, draw=ln-rule] (21,-11.8) rectangle ++(1,-1);
  \node[anchor=west, text=ln-muted] at (22.3,-12.3)
    {\iflangja{ミス待ち（ストール）}{waiting for the miss (stall)}};
  \node[anchor=west, text=ln-muted] at (22.3,-13.5)
    {\iflangja{I4 だけが I1 の R1 に依存する}{only I4 depends on R1 from I1}};
\end{tikzpicture}

  \caption{One program whose first instruction misses in the cache:
    (a) in-order, every later instruction waits; (b) out-of-order, only
    the dependent \texttt{I4} waits; (c) multithreading, four threads
    each issue their load and the four misses overlap.}
  \label{fig:l04-latency-hiding}
\end{figure}

\begin{definition}[Latency hiding by multithreading]
  \label{def:l04-latency-hiding}
  Keeping the pipeline busy with useful instructions while a
  long-latency operation (a cache miss, a long ALU operation, a branch
  resolution) is in flight is called
  \term{latency hiding}[レイテンシ隠蔽].  Holding the state of many
  threads on chip and switching among them at every cycle, so that a
  stalled thread simply is not selected, is
  \term{fine-grained multithreading}[細粒度マルチスレッディング].  The number of memory requests a processor keeps in flight
  at once is its \term{memory-level parallelism (MLP)}[メモリレベル並列性]\index{MLP|see{memory-level parallelism (MLP)}};
  the multithreaded pipeline of \cref{fig:l04-latency-hiding}c has an
  MLP of four.
\end{definition}

Where a CPU spends its area on out-of-order logic, large caches and the
search for ILP, a GPU spends it on threads: the stalls caused by cache
misses, branches and long-latency arithmetic are covered by the other
warps of the SM.  The price is stated by one rule: the longer the memory
latency, the more threads are needed to hide it.\tip{Threads needed
$\approx$ latency $\div$ cycles of independent work each thread
supplies between two long-latency operations.  Doubling the latency
doubles the thread count; so does halving the independent work.}

\begin{example}
  \label{ex:l04-threads-needed}
  An SM issues one warp instruction per cycle.  Each warp, after issuing
  a load, can issue two more instructions before it needs the loaded
  value, and the load takes \num{400} cycles.  To keep the issue slot
  busy for the whole \num{400} cycles the SM needs about
  $400 / 3 \approx 134$ warps that can each contribute three
  instructions; with a \num{800}-cycle latency it needs about
  \num{267}.  No real SM holds that many warps (the limit is
  \num{64}), which is why GPUs also rely on caches, on independent
  instructions within a warp and, as \cref{sec:l04-coalescing} shows,
  on cutting the number of memory requests.\margin{Maximum resident
  warps per SM: 64 on compute capabilities 7.0, 8.0 and 9.0, 48 on 8.6
  and 8.9, 32 on 7.5: CPG Table~21 \cite{nvidia2024cudaguide}; the
  occupancy of Lesson~3 is measured against this limit.}
\end{example}

\subsection{Program counters and registers per thread}

Multithreading changes the front end of the pipeline
(\cref{fig:l04-warp-state}).  A single-threaded core keeps one
\term{program counter (PC)}[プログラムカウンタ]\index{PC|see{program counter (PC)}}, the address of the next
instruction to fetch.  An SM keeps one PC per warp: four resident warps
mean four PC registers, and in every cycle the fetch stage reads the PC
of the warp it has selected and fetches one instruction for all the
threads of that warp.  One static instruction thus serves a whole warp,
which is the SPMD model of Lesson~1 realised in hardware.  Behind the
front end, the \term{register file}[レジスタファイル] holds a private set
of registers for every thread of every resident warp, so that switching
from one warp to another needs no copying: the registers of the next
warp are already there.\margin{Per SM the register file holds
\num{65536} 32-bit registers, \SI{256}{\kibi\byte}, on every compute
capability from 5.0 to 9.0: CPG Table~21 \cite{nvidia2024cudaguide}.
Fermi (2010) had \SI{128}{\kibi\byte} per SM: K12 §1.4.1 \cite{kim2012gpgpu}.}

\begin{figure}[htbp]
  \centering
  % Front-end state of a multithreaded SM (Lesson 4): one program counter
% per warp (a PC table W0..W3), feeding the instruction cache and decoder,
% and a register file with a private register set for every thread,
% grouped by warp.  Redraws the outline figure m4l1_fe_extension.
% Usage: % Front-end state of a multithreaded SM (Lesson 4): one program counter
% per warp (a PC table W0..W3), feeding the instruction cache and decoder,
% and a register file with a private register set for every thread,
% grouped by warp.  Redraws the outline figure m4l1_fe_extension.
% Usage: % Front-end state of a multithreaded SM (Lesson 4): one program counter
% per warp (a PC table W0..W3), feeding the instruction cache and decoder,
% and a register file with a private register set for every thread,
% grouped by warp.  Redraws the outline figure m4l1_fe_extension.
% Usage: \input{figures/l04-warp-state}
% Width: body (120 mm); the drawing is about 118 mm wide.
\begin{tikzpicture}[x=1mm, y=1mm, font=\scriptsize\sffamily,
    >={Stealth[length=1.6mm]},
    blk/.style={draw, fill=ln-box, minimum height=6mm, inner sep=2pt,
                align=center}]
  % PC table
  \node[anchor=south west, font=\scriptsize\sffamily\bfseries] at (0,27)
    {\iflangja{ワープごとの PC}{one PC per warp}};
  \foreach \w/\pc [count=\i from 0] in {W0/0x8000, W1/0x8004, W2/0x8000, W3/0x800C} {
    \draw[fill=ln-accent!12] (0,26-\i*6) rectangle ++(8,-5.5)
      node[midway] {\w};
    \draw[fill=white] (8,26-\i*6) rectangle ++(14,-5.5)
      node[midway, font=\scriptsize\ttfamily] {\pc};
  }
  % fetch path
  \node[blk, minimum width=14mm] (icache) at (36,14)
    {\iflangja{命令\\キャッシュ}{I-cache}};
  \node[blk, minimum width=14mm, right=3mm of icache] (dec)
    {\iflangja{デコーダ}{Decoder}};
  \draw[->] (22,14) -- (icache);
  \draw[->] (icache) -- (dec);
  \node[anchor=north, text=ln-muted, align=center] at (36,8.5)
    {\iflangja{選ばれたワープの\\PC で1命令をフェッチ}{one instruction fetched\\at the selected warp's PC}};
  % register file: 4 warps x 4 threads, registers R0..R3
  \node[anchor=south west, font=\scriptsize\sffamily\bfseries] at (66,27)
    {\iflangja{スレッドごとのレジスタ}{a register set per thread}};
  \foreach \w in {0,...,3} {
    \draw[fill=ln-accent!12] (66+\w*13,26) rectangle ++(12.6,-4)
      node[midway] {W\w};
    \foreach \t in {0,...,3} {
      \node[font=\scriptsize\sffamily] at (66+\w*13+\t*3.2+1.6,20.2) {T\t};
    }
  }
  \foreach \r in {0,...,3} {
    \node[anchor=east, inner sep=1pt] at (65.5,16.6-\r*3.4) {R\r};
    \foreach \w in {0,...,3} {
      \foreach \t in {0,...,3} {
        \draw[fill=white] (66+\w*13+\t*3.2,18.3-\r*3.4) rectangle ++(3.2,-3.4);
      }
    }
  }
  \node[anchor=north west, text=ln-muted, align=left] at (66,3)
    {\iflangja{W0 の R1 読み出し = T0--T3 の R1 を同時に読む}{reading R1 for W0 = R1 of T0--T3 at once}};
\end{tikzpicture}

% Width: body (120 mm); the drawing is about 118 mm wide.
\begin{tikzpicture}[x=1mm, y=1mm, font=\scriptsize\sffamily,
    >={Stealth[length=1.6mm]},
    blk/.style={draw, fill=ln-box, minimum height=6mm, inner sep=2pt,
                align=center}]
  % PC table
  \node[anchor=south west, font=\scriptsize\sffamily\bfseries] at (0,27)
    {\iflangja{ワープごとの PC}{one PC per warp}};
  \foreach \w/\pc [count=\i from 0] in {W0/0x8000, W1/0x8004, W2/0x8000, W3/0x800C} {
    \draw[fill=ln-accent!12] (0,26-\i*6) rectangle ++(8,-5.5)
      node[midway] {\w};
    \draw[fill=white] (8,26-\i*6) rectangle ++(14,-5.5)
      node[midway, font=\scriptsize\ttfamily] {\pc};
  }
  % fetch path
  \node[blk, minimum width=14mm] (icache) at (36,14)
    {\iflangja{命令\\キャッシュ}{I-cache}};
  \node[blk, minimum width=14mm, right=3mm of icache] (dec)
    {\iflangja{デコーダ}{Decoder}};
  \draw[->] (22,14) -- (icache);
  \draw[->] (icache) -- (dec);
  \node[anchor=north, text=ln-muted, align=center] at (36,8.5)
    {\iflangja{選ばれたワープの\\PC で1命令をフェッチ}{one instruction fetched\\at the selected warp's PC}};
  % register file: 4 warps x 4 threads, registers R0..R3
  \node[anchor=south west, font=\scriptsize\sffamily\bfseries] at (66,27)
    {\iflangja{スレッドごとのレジスタ}{a register set per thread}};
  \foreach \w in {0,...,3} {
    \draw[fill=ln-accent!12] (66+\w*13,26) rectangle ++(12.6,-4)
      node[midway] {W\w};
    \foreach \t in {0,...,3} {
      \node[font=\scriptsize\sffamily] at (66+\w*13+\t*3.2+1.6,20.2) {T\t};
    }
  }
  \foreach \r in {0,...,3} {
    \node[anchor=east, inner sep=1pt] at (65.5,16.6-\r*3.4) {R\r};
    \foreach \w in {0,...,3} {
      \foreach \t in {0,...,3} {
        \draw[fill=white] (66+\w*13+\t*3.2,18.3-\r*3.4) rectangle ++(3.2,-3.4);
      }
    }
  }
  \node[anchor=north west, text=ln-muted, align=left] at (66,3)
    {\iflangja{W0 の R1 読み出し = T0--T3 の R1 を同時に読む}{reading R1 for W0 = R1 of T0--T3 at once}};
\end{tikzpicture}

% Width: body (120 mm); the drawing is about 118 mm wide.
\begin{tikzpicture}[x=1mm, y=1mm, font=\scriptsize\sffamily,
    >={Stealth[length=1.6mm]},
    blk/.style={draw, fill=ln-box, minimum height=6mm, inner sep=2pt,
                align=center}]
  % PC table
  \node[anchor=south west, font=\scriptsize\sffamily\bfseries] at (0,27)
    {\iflangja{ワープごとの PC}{one PC per warp}};
  \foreach \w/\pc [count=\i from 0] in {W0/0x8000, W1/0x8004, W2/0x8000, W3/0x800C} {
    \draw[fill=ln-accent!12] (0,26-\i*6) rectangle ++(8,-5.5)
      node[midway] {\w};
    \draw[fill=white] (8,26-\i*6) rectangle ++(14,-5.5)
      node[midway, font=\scriptsize\ttfamily] {\pc};
  }
  % fetch path
  \node[blk, minimum width=14mm] (icache) at (36,14)
    {\iflangja{命令\\キャッシュ}{I-cache}};
  \node[blk, minimum width=14mm, right=3mm of icache] (dec)
    {\iflangja{デコーダ}{Decoder}};
  \draw[->] (22,14) -- (icache);
  \draw[->] (icache) -- (dec);
  \node[anchor=north, text=ln-muted, align=center] at (36,8.5)
    {\iflangja{選ばれたワープの\\PC で1命令をフェッチ}{one instruction fetched\\at the selected warp's PC}};
  % register file: 4 warps x 4 threads, registers R0..R3
  \node[anchor=south west, font=\scriptsize\sffamily\bfseries] at (66,27)
    {\iflangja{スレッドごとのレジスタ}{a register set per thread}};
  \foreach \w in {0,...,3} {
    \draw[fill=ln-accent!12] (66+\w*13,26) rectangle ++(12.6,-4)
      node[midway] {W\w};
    \foreach \t in {0,...,3} {
      \node[font=\scriptsize\sffamily] at (66+\w*13+\t*3.2+1.6,20.2) {T\t};
    }
  }
  \foreach \r in {0,...,3} {
    \node[anchor=east, inner sep=1pt] at (65.5,16.6-\r*3.4) {R\r};
    \foreach \w in {0,...,3} {
      \foreach \t in {0,...,3} {
        \draw[fill=white] (66+\w*13+\t*3.2,18.3-\r*3.4) rectangle ++(3.2,-3.4);
      }
    }
  }
  \node[anchor=north west, text=ln-muted, align=left] at (66,3)
    {\iflangja{W0 の R1 読み出し = T0--T3 の R1 を同時に読む}{reading R1 for W0 = R1 of T0--T3 at once}};
\end{tikzpicture}

  \caption{The front end of a multithreaded SM: one program counter per
    warp selects the instruction to fetch, and the register file keeps a
    register set for every thread, grouped by warp, so that reading
    \texttt{R1} for warp~0 reads \texttt{R1} of its four threads at
    once.}
  \label{fig:l04-warp-state}
\end{figure}

\begin{definition}[Context switch]
  \label{def:l04-context-switch}
  A \term{context switch}[コンテキストスイッチ] replaces the thread whose
  state (PC, registers) the pipeline is executing with another thread.
  On a CPU the pipeline, the PC and the registers belong to one thread
  at a time, so the switch stores the running thread's state to memory
  and loads the next thread's (\cref{fig:l04-context-switch}), a cost
  of many cycles.  On a GPU every resident warp's PC and registers are
  already on chip; the switch only changes which PC and which register
  set the pipeline uses, and costs nothing.\margin{CPG §4.2 Hardware
  Multithreading: the execution context of each warp \enquote{is
  maintained on-chip during the entire lifetime of the warp.  Therefore,
  switching from one execution context to another has no cost}
  \cite{nvidia2024cudaguide}.}
\end{definition}

\begin{figure}[htbp]
  \centering
  % A CPU context switch (Lesson 4): the core holds one thread's PC and
% registers; switching from T1 to T2 stores T1's state to memory and loads
% T2's.  Redraws the outline figures m4l1_cpu_context, _1 and _2.
% Usage: % A CPU context switch (Lesson 4): the core holds one thread's PC and
% registers; switching from T1 to T2 stores T1's state to memory and loads
% T2's.  Redraws the outline figures m4l1_cpu_context, _1 and _2.
% Usage: % A CPU context switch (Lesson 4): the core holds one thread's PC and
% registers; switching from T1 to T2 stores T1's state to memory and loads
% T2's.  Redraws the outline figures m4l1_cpu_context, _1 and _2.
% Usage: \input{figures/l04-context-switch}
% Width: body (120 mm); the drawing is about 100 mm wide.
\begin{tikzpicture}[x=1mm, y=1mm, font=\scriptsize\sffamily,
    >={Stealth[length=1.6mm]},
    blk/.style={draw, fill=ln-box, minimum height=5mm, inner sep=2pt,
                align=center}]
  % CPU core
  \draw[ln-accent, thick, rounded corners=1pt] (0,0) rectangle (44,34);
  \node[anchor=north west, font=\scriptsize\sffamily\bfseries, text=ln-accent]
    at (1,34) {\iflangja{CPU コア（1スレッド分の状態）}{CPU core (state of one thread)}};
  \node[blk, anchor=west, minimum width=12mm] at (2,25) {\iflangja{命令\\キャッシュ}{I-cache}};
  \node[blk, anchor=west, minimum width=12mm] at (16,25) {\iflangja{デコーダ}{Decoder}};
  \node[blk, anchor=west, minimum width=12mm] at (30,25) {\iflangja{スケジューラ}{Sched.}};
  \foreach \k in {0,...,3} {
    \draw[fill=ln-accent!15] (2+\k*7,14) rectangle ++(6,-5) node[midway] {ALU};
  }
  % one register set, coloured for the resident thread
  \node[anchor=south west, inner sep=1pt] at (31,16.5) {\iflangja{レジスタ}{registers}};
  \foreach \r/\lbl [count=\i from 0] in {PC/PC, R0/R0, R1/R1, R2/R2} {
    \draw[fill=ln-tip!25] (31,16-\i*3) rectangle ++(11,-3) node[midway] {\lbl};
  }
  \node[anchor=north west, text=ln-tip, font=\scriptsize\sffamily\bfseries] at (30,3.5)
    {\iflangja{= T1 の値}{= T1's values}};
  % memory with saved contexts
  \draw[ln-accent, thick, rounded corners=1pt] (70,0) rectangle (100,34);
  \node[anchor=north west, font=\scriptsize\sffamily\bfseries, text=ln-accent]
    at (71,34) {\iflangja{メモリ}{Memory}};
  \foreach \t/\c [count=\i from 0] in {T0/ln-accent!25, T1/ln-tip!25, T2/ln-caution!20, T3/ln-source!20} {
    \draw[fill=\c] (72,28-\i*7) rectangle ++(26,-6)
      node[midway] {\iflangja{\t のコンテキスト}{context of \t}};
  }
  % switch arrows
  \draw[->, thick, ln-tip] (42,10) -- (72,18.5)
    node[pos=0.45, above, sloped, inner sep=1pt] {\iflangja{1. T1 を退避}{1. save T1}};
  \draw[->, thick, ln-caution] (72,11) -- (42,5)
    node[pos=0.55, below, sloped, inner sep=1pt] {\iflangja{2. T2 を復元}{2. restore T2}};
\end{tikzpicture}

% Width: body (120 mm); the drawing is about 100 mm wide.
\begin{tikzpicture}[x=1mm, y=1mm, font=\scriptsize\sffamily,
    >={Stealth[length=1.6mm]},
    blk/.style={draw, fill=ln-box, minimum height=5mm, inner sep=2pt,
                align=center}]
  % CPU core
  \draw[ln-accent, thick, rounded corners=1pt] (0,0) rectangle (44,34);
  \node[anchor=north west, font=\scriptsize\sffamily\bfseries, text=ln-accent]
    at (1,34) {\iflangja{CPU コア（1スレッド分の状態）}{CPU core (state of one thread)}};
  \node[blk, anchor=west, minimum width=12mm] at (2,25) {\iflangja{命令\\キャッシュ}{I-cache}};
  \node[blk, anchor=west, minimum width=12mm] at (16,25) {\iflangja{デコーダ}{Decoder}};
  \node[blk, anchor=west, minimum width=12mm] at (30,25) {\iflangja{スケジューラ}{Sched.}};
  \foreach \k in {0,...,3} {
    \draw[fill=ln-accent!15] (2+\k*7,14) rectangle ++(6,-5) node[midway] {ALU};
  }
  % one register set, coloured for the resident thread
  \node[anchor=south west, inner sep=1pt] at (31,16.5) {\iflangja{レジスタ}{registers}};
  \foreach \r/\lbl [count=\i from 0] in {PC/PC, R0/R0, R1/R1, R2/R2} {
    \draw[fill=ln-tip!25] (31,16-\i*3) rectangle ++(11,-3) node[midway] {\lbl};
  }
  \node[anchor=north west, text=ln-tip, font=\scriptsize\sffamily\bfseries] at (30,3.5)
    {\iflangja{= T1 の値}{= T1's values}};
  % memory with saved contexts
  \draw[ln-accent, thick, rounded corners=1pt] (70,0) rectangle (100,34);
  \node[anchor=north west, font=\scriptsize\sffamily\bfseries, text=ln-accent]
    at (71,34) {\iflangja{メモリ}{Memory}};
  \foreach \t/\c [count=\i from 0] in {T0/ln-accent!25, T1/ln-tip!25, T2/ln-caution!20, T3/ln-source!20} {
    \draw[fill=\c] (72,28-\i*7) rectangle ++(26,-6)
      node[midway] {\iflangja{\t のコンテキスト}{context of \t}};
  }
  % switch arrows
  \draw[->, thick, ln-tip] (42,10) -- (72,18.5)
    node[pos=0.45, above, sloped, inner sep=1pt] {\iflangja{1. T1 を退避}{1. save T1}};
  \draw[->, thick, ln-caution] (72,11) -- (42,5)
    node[pos=0.55, below, sloped, inner sep=1pt] {\iflangja{2. T2 を復元}{2. restore T2}};
\end{tikzpicture}

% Width: body (120 mm); the drawing is about 100 mm wide.
\begin{tikzpicture}[x=1mm, y=1mm, font=\scriptsize\sffamily,
    >={Stealth[length=1.6mm]},
    blk/.style={draw, fill=ln-box, minimum height=5mm, inner sep=2pt,
                align=center}]
  % CPU core
  \draw[ln-accent, thick, rounded corners=1pt] (0,0) rectangle (44,34);
  \node[anchor=north west, font=\scriptsize\sffamily\bfseries, text=ln-accent]
    at (1,34) {\iflangja{CPU コア（1スレッド分の状態）}{CPU core (state of one thread)}};
  \node[blk, anchor=west, minimum width=12mm] at (2,25) {\iflangja{命令\\キャッシュ}{I-cache}};
  \node[blk, anchor=west, minimum width=12mm] at (16,25) {\iflangja{デコーダ}{Decoder}};
  \node[blk, anchor=west, minimum width=12mm] at (30,25) {\iflangja{スケジューラ}{Sched.}};
  \foreach \k in {0,...,3} {
    \draw[fill=ln-accent!15] (2+\k*7,14) rectangle ++(6,-5) node[midway] {ALU};
  }
  % one register set, coloured for the resident thread
  \node[anchor=south west, inner sep=1pt] at (31,16.5) {\iflangja{レジスタ}{registers}};
  \foreach \r/\lbl [count=\i from 0] in {PC/PC, R0/R0, R1/R1, R2/R2} {
    \draw[fill=ln-tip!25] (31,16-\i*3) rectangle ++(11,-3) node[midway] {\lbl};
  }
  \node[anchor=north west, text=ln-tip, font=\scriptsize\sffamily\bfseries] at (30,3.5)
    {\iflangja{= T1 の値}{= T1's values}};
  % memory with saved contexts
  \draw[ln-accent, thick, rounded corners=1pt] (70,0) rectangle (100,34);
  \node[anchor=north west, font=\scriptsize\sffamily\bfseries, text=ln-accent]
    at (71,34) {\iflangja{メモリ}{Memory}};
  \foreach \t/\c [count=\i from 0] in {T0/ln-accent!25, T1/ln-tip!25, T2/ln-caution!20, T3/ln-source!20} {
    \draw[fill=\c] (72,28-\i*7) rectangle ++(26,-6)
      node[midway] {\iflangja{\t のコンテキスト}{context of \t}};
  }
  % switch arrows
  \draw[->, thick, ln-tip] (42,10) -- (72,18.5)
    node[pos=0.45, above, sloped, inner sep=1pt] {\iflangja{1. T1 を退避}{1. save T1}};
  \draw[->, thick, ln-caution] (72,11) -- (42,5)
    node[pos=0.55, below, sloped, inner sep=1pt] {\iflangja{2. T2 を復元}{2. restore T2}};
\end{tikzpicture}

  \caption{A CPU context switch: the core holds one thread's PC and
    registers, so switching from \texttt{T1} to \texttt{T2} saves
    \texttt{T1}'s state to memory and restores \texttt{T2}'s.}
  \label{fig:l04-context-switch}
\end{figure}

\begin{note}
  The word \emph{context switch} is used in two senses.  An operating
  system's context switch saves a process or thread to memory and is
  expensive; the GPU's \enquote{context switch} between
  warps is a selection among states that are all resident and is
  free.\caution{The free GPU switch applies only to warps that are
  resident on the SM.  Replacing a finished thread block by a new one
  does allocate registers and shared memory and is not free; Lesson~7
  treats warp scheduling.}  A CPU with simultaneous
  multithreading holds two thread contexts per core; an SM holds
  hundreds of threads.\margin{K12 §1.2.1 contrasts the two hyper-threads
  of a CPU core with the hundreds of threads of an SM
  \cite{kim2012gpgpu}.}
\end{note}

\subsection{Resource requirements}

The state that multithreading keeps on chip is counted in two
quantities.
\begin{itemize}
  \item \emph{Program counters.}  One per warp, so the number of PCs is
    the number of resident threads divided by the warp size.  Later
    GPUs keep more front-end state than this; for the model of this
    book one PC per warp is assumed.\margin{Since Volta, independent
    thread scheduling keeps a program counter and call stack per
    \emph{thread}, so that diverged threads can be interleaved: CPG §4.1
    SIMT Architecture \cite{nvidia2024cudaguide}.  Lesson~5 discusses
    it; Lesson~6 keeps the per-warp model.}
  \item \emph{Registers.}  $K$ registers per thread times the number of
    resident threads, where $K$ is chosen by the compiler for each
    kernel and varies from application to application.  A larger $K$
    means fewer resident threads for a fixed register file, which is the
    register limit of the occupancy calculation of Lesson~3.
\end{itemize}

\begin{example}
  \label{ex:l04-resources}
  An SM holds \num{512} resident threads with a warp size of 32, so it
  keeps $512 / 32 = 16$ program counters.  A kernel of \num{12}
  instructions executed by all of them costs
  $16 \times 12 = 192$ warp-instruction fetches, whereas the same work
  on single-threaded cores would need $512 \times 12 = \num{6144}$
  instruction fetches: one fetch serves 32 threads.  If the compiler
  assigns $K = 40$ registers per thread, the resident threads occupy
  $40 \times 512 = \num{20480}$ registers, \SI{80}{\kibi\byte} at
  \SI{4}{\byte} each.\tip{Per SM: PCs = threads $\div$ warp size;
  registers = $K \times$ threads; a GPU repeats both in every SM.  Both scale with the thread count, so the
  latency-hiding rule above is a statement about on-chip storage.}
\end{example}

\begin{keypoint}
  A GPU hides memory latency with fine-grained multithreading rather
  than with out-of-order execution: every resident warp's PC and
  registers stay on chip, so switching warps costs nothing, and the
  number of warps an SM can hold is a storage question (one PC per
  warp, $K$ registers per thread).
\end{keypoint}

\section{Bank conflicts}
\label{sec:l04-bank-conflicts}

The register file of the previous section must deliver operands for a
whole warp at once, which makes its bandwidth, not its capacity, the
design problem.\source{M4, bank conflicts; K12 §1.4.1 (register
read/write, scoreboarding); A18 ch.~3 (operand collector); BPG §10.2.3
Shared Memory; CPG Compute Capability 5.x, Shared Memory;
\cite{harris2013sharedmem}.}  This section introduces ports and banks,
the conflicts banking creates in the register file and in shared memory,
and the operand collector and scoreboard that absorb them.

\subsection{Register file bandwidth}

A typical arithmetic instruction reads two source registers and writes
one destination.  An SM executes an instruction for many threads at
once: with an execution width of 16 lanes, one instruction moves
$16 \times 3 = 48$ register values in a cycle, and a wider SM
proportionally more.\margin{K12 §1.4.1 scales this up: a GPU of
\SI{1}{\tera\flop\per\second} with two reads and one write per
operation needs on the order of \SI{64}{\tera\byte\per\second} of
register read bandwidth \cite{kim2012gpgpu}, far beyond any single
memory array.}  A memory array that can be read and written through
that many paths at once is what the hardware must provide.

\begin{definition}[Port]
  \label{def:l04-port}
  A \term{port}[ポート] is a hardware interface through which one value
  can be read from or written to a storage array in a cycle; a
  \term{read port}[読み出しポート] carries a value out, a
  \term{write port}[書き込みポート] carries one in.  Each port is a set
  of wires and selection logic attached to every entry of the array,
  so its cost grows with the number of entries, and an array with many
  ports is large and slow.
\end{definition}

\Cref{fig:l04-ports-vs-banks}a shows the brute-force design for an
execution width of 4: every register entry carries $2 \times 4 = 8$
read ports and 4 write ports, so that any 8 values can be read and any
4 written in one cycle.  \Cref{fig:l04-ports-vs-banks}b reaches the
same throughput more cheaply by splitting the array.

\begin{definition}[Bank]
  \label{def:l04-bank}
  A \term{bank}[バンク] is one of several independent storage arrays
  into which a register file or memory is divided, each with its own
  (few) ports.  Banks are accessed simultaneously, so $B$ banks with
  $p$ ports each deliver up to $Bp$ values per cycle as long as the
  values requested lie in different banks.  In
  \cref{fig:l04-ports-vs-banks}b four banks with 2 read and 1 write
  port each give the 8 reads and 4 writes of the unbanked design with
  a quarter of the ports per entry.
\end{definition}

\begin{figure}[htbp]
  \centering
  % Ports versus banks in a register file for an execution width of 4
% (Lesson 4).  Left: without banks every register entry needs 8 read and
% 4 write ports.  Right: four banks, each with 2 read and 1 write port,
% accessed simultaneously.  Redraws the outline figures m4l2_port_vs_bank
% and m4l2_port_vs_bank_1.
% Usage: % Ports versus banks in a register file for an execution width of 4
% (Lesson 4).  Left: without banks every register entry needs 8 read and
% 4 write ports.  Right: four banks, each with 2 read and 1 write port,
% accessed simultaneously.  Redraws the outline figures m4l2_port_vs_bank
% and m4l2_port_vs_bank_1.
% Usage: % Ports versus banks in a register file for an execution width of 4
% (Lesson 4).  Left: without banks every register entry needs 8 read and
% 4 write ports.  Right: four banks, each with 2 read and 1 write port,
% accessed simultaneously.  Redraws the outline figures m4l2_port_vs_bank
% and m4l2_port_vs_bank_1.
% Usage: \input{figures/l04-ports-vs-banks}
% Width: body (120 mm); the drawing is about 118 mm wide.
\begin{tikzpicture}[x=1mm, y=1mm, font=\scriptsize\sffamily]
  % (a) no banks: one column of entries crossed by 8 R and 4 W lines
  \node[anchor=south west, font=\scriptsize\sffamily\bfseries] at (0,43)
    {\iflangja{(a) バンクなし}{(a) no banks}};
  \foreach \t in {0,1} \foreach \r in {0,...,3} {
    \pgfmathsetmacro{\yy}{38-(\t*4+\r)*4.2}
    \draw[fill=ln-box] (0,\yy) rectangle ++(12,-4) node[midway] {T\t.R\r};
  }
  \node[anchor=north] at (6,3) {$\vdots$};
  \foreach \k in {0,...,7} {
    \draw[ln-accent, thick] (14+\k*2.2,39.5) -- ++(0,-35.5);
    \node[text=ln-accent, anchor=south] at (14+\k*2.2,39.6) {R};
  }
  \foreach \k in {0,...,3} {
    \draw[ln-caution, thick] (33+\k*2.2,39.5) -- ++(0,-35.5);
    \node[text=ln-caution, anchor=south] at (33+\k*2.2,39.6) {W};
  }
  \foreach \t in {0,1} \foreach \r in {0,...,3} {
    \pgfmathsetmacro{\yy}{36-(\t*4+\r)*4.2}
    \draw[ln-muted] (12,\yy) -- (42,\yy);
  }
  \node[anchor=north west, align=left, text=ln-muted] at (0,-1)
    {\iflangja{各エントリに 8 読み出し + 4 書き込みポート}{8 read + 4 write ports on every entry}};
  % (b) four banks
  \node[anchor=south west, font=\scriptsize\sffamily\bfseries] at (58,43)
    {\iflangja{(b) 4 バンク}{(b) 4 banks}};
  \foreach \b in {0,...,3} {
    \pgfmathsetmacro{\xx}{58+\b*15}
    \node[anchor=south, text=ln-muted] at (\xx+6,21.5) {\iflangja{バンク \b}{bank \b}};
    \foreach \r in {0,...,3} {
      \pgfmathsetmacro{\yy}{38-\r*4.2}
      \draw[fill=ln-box] (\xx,\yy) rectangle ++(9,-4) node[midway] {T\b.R\r};
      \draw[ln-muted] (\xx+9,\yy-2) -- (\xx+14.5,\yy-2);
    }
    \draw[ln-accent, thick] (\xx+10.2,39.5) -- ++(0,-18);
    \draw[ln-accent, thick] (\xx+12.2,39.5) -- ++(0,-18);
    \draw[ln-caution, thick] (\xx+14.2,39.5) -- ++(0,-18);
    \node[text=ln-accent, anchor=south] at (\xx+10.2,39.6) {R};
    \node[text=ln-accent, anchor=south] at (\xx+12.2,39.6) {R};
    \node[text=ln-caution, anchor=south] at (\xx+14.2,39.6) {W};
  }
  \node[anchor=north west, align=left, text=ln-muted] at (58,17)
    {\iflangja{各バンクに 2 読み出し + 1 書き込みポート；\\4 バンクを同時にアクセスできる}{2 read + 1 write ports per bank;\\the four banks are accessed at the same time}};
\end{tikzpicture}

% Width: body (120 mm); the drawing is about 118 mm wide.
\begin{tikzpicture}[x=1mm, y=1mm, font=\scriptsize\sffamily]
  % (a) no banks: one column of entries crossed by 8 R and 4 W lines
  \node[anchor=south west, font=\scriptsize\sffamily\bfseries] at (0,43)
    {\iflangja{(a) バンクなし}{(a) no banks}};
  \foreach \t in {0,1} \foreach \r in {0,...,3} {
    \pgfmathsetmacro{\yy}{38-(\t*4+\r)*4.2}
    \draw[fill=ln-box] (0,\yy) rectangle ++(12,-4) node[midway] {T\t.R\r};
  }
  \node[anchor=north] at (6,3) {$\vdots$};
  \foreach \k in {0,...,7} {
    \draw[ln-accent, thick] (14+\k*2.2,39.5) -- ++(0,-35.5);
    \node[text=ln-accent, anchor=south] at (14+\k*2.2,39.6) {R};
  }
  \foreach \k in {0,...,3} {
    \draw[ln-caution, thick] (33+\k*2.2,39.5) -- ++(0,-35.5);
    \node[text=ln-caution, anchor=south] at (33+\k*2.2,39.6) {W};
  }
  \foreach \t in {0,1} \foreach \r in {0,...,3} {
    \pgfmathsetmacro{\yy}{36-(\t*4+\r)*4.2}
    \draw[ln-muted] (12,\yy) -- (42,\yy);
  }
  \node[anchor=north west, align=left, text=ln-muted] at (0,-1)
    {\iflangja{各エントリに 8 読み出し + 4 書き込みポート}{8 read + 4 write ports on every entry}};
  % (b) four banks
  \node[anchor=south west, font=\scriptsize\sffamily\bfseries] at (58,43)
    {\iflangja{(b) 4 バンク}{(b) 4 banks}};
  \foreach \b in {0,...,3} {
    \pgfmathsetmacro{\xx}{58+\b*15}
    \node[anchor=south, text=ln-muted] at (\xx+6,21.5) {\iflangja{バンク \b}{bank \b}};
    \foreach \r in {0,...,3} {
      \pgfmathsetmacro{\yy}{38-\r*4.2}
      \draw[fill=ln-box] (\xx,\yy) rectangle ++(9,-4) node[midway] {T\b.R\r};
      \draw[ln-muted] (\xx+9,\yy-2) -- (\xx+14.5,\yy-2);
    }
    \draw[ln-accent, thick] (\xx+10.2,39.5) -- ++(0,-18);
    \draw[ln-accent, thick] (\xx+12.2,39.5) -- ++(0,-18);
    \draw[ln-caution, thick] (\xx+14.2,39.5) -- ++(0,-18);
    \node[text=ln-accent, anchor=south] at (\xx+10.2,39.6) {R};
    \node[text=ln-accent, anchor=south] at (\xx+12.2,39.6) {R};
    \node[text=ln-caution, anchor=south] at (\xx+14.2,39.6) {W};
  }
  \node[anchor=north west, align=left, text=ln-muted] at (58,17)
    {\iflangja{各バンクに 2 読み出し + 1 書き込みポート；\\4 バンクを同時にアクセスできる}{2 read + 1 write ports per bank;\\the four banks are accessed at the same time}};
\end{tikzpicture}

% Width: body (120 mm); the drawing is about 118 mm wide.
\begin{tikzpicture}[x=1mm, y=1mm, font=\scriptsize\sffamily]
  % (a) no banks: one column of entries crossed by 8 R and 4 W lines
  \node[anchor=south west, font=\scriptsize\sffamily\bfseries] at (0,43)
    {\iflangja{(a) バンクなし}{(a) no banks}};
  \foreach \t in {0,1} \foreach \r in {0,...,3} {
    \pgfmathsetmacro{\yy}{38-(\t*4+\r)*4.2}
    \draw[fill=ln-box] (0,\yy) rectangle ++(12,-4) node[midway] {T\t.R\r};
  }
  \node[anchor=north] at (6,3) {$\vdots$};
  \foreach \k in {0,...,7} {
    \draw[ln-accent, thick] (14+\k*2.2,39.5) -- ++(0,-35.5);
    \node[text=ln-accent, anchor=south] at (14+\k*2.2,39.6) {R};
  }
  \foreach \k in {0,...,3} {
    \draw[ln-caution, thick] (33+\k*2.2,39.5) -- ++(0,-35.5);
    \node[text=ln-caution, anchor=south] at (33+\k*2.2,39.6) {W};
  }
  \foreach \t in {0,1} \foreach \r in {0,...,3} {
    \pgfmathsetmacro{\yy}{36-(\t*4+\r)*4.2}
    \draw[ln-muted] (12,\yy) -- (42,\yy);
  }
  \node[anchor=north west, align=left, text=ln-muted] at (0,-1)
    {\iflangja{各エントリに 8 読み出し + 4 書き込みポート}{8 read + 4 write ports on every entry}};
  % (b) four banks
  \node[anchor=south west, font=\scriptsize\sffamily\bfseries] at (58,43)
    {\iflangja{(b) 4 バンク}{(b) 4 banks}};
  \foreach \b in {0,...,3} {
    \pgfmathsetmacro{\xx}{58+\b*15}
    \node[anchor=south, text=ln-muted] at (\xx+6,21.5) {\iflangja{バンク \b}{bank \b}};
    \foreach \r in {0,...,3} {
      \pgfmathsetmacro{\yy}{38-\r*4.2}
      \draw[fill=ln-box] (\xx,\yy) rectangle ++(9,-4) node[midway] {T\b.R\r};
      \draw[ln-muted] (\xx+9,\yy-2) -- (\xx+14.5,\yy-2);
    }
    \draw[ln-accent, thick] (\xx+10.2,39.5) -- ++(0,-18);
    \draw[ln-accent, thick] (\xx+12.2,39.5) -- ++(0,-18);
    \draw[ln-caution, thick] (\xx+14.2,39.5) -- ++(0,-18);
    \node[text=ln-accent, anchor=south] at (\xx+10.2,39.6) {R};
    \node[text=ln-accent, anchor=south] at (\xx+12.2,39.6) {R};
    \node[text=ln-caution, anchor=south] at (\xx+14.2,39.6) {W};
  }
  \node[anchor=north west, align=left, text=ln-muted] at (58,17)
    {\iflangja{各バンクに 2 読み出し + 1 書き込みポート；\\4 バンクを同時にアクセスできる}{2 read + 1 write ports per bank;\\the four banks are accessed at the same time}};
\end{tikzpicture}

  \caption{Ports versus banks for an execution width of 4.  (a) Without
    banks every entry needs 8 read and 4 write ports.  (b) Four banks
    with 2 read and 1 write ports each are accessed at the same time.}
  \label{fig:l04-ports-vs-banks}
\end{figure}

\subsection{Register bank conflicts}

The saving has a condition: the values a cycle needs must be spread over
the banks.

\begin{definition}[Bank conflict]
  \label{def:l04-bank-conflict}
  A \term{bank conflict}[バンク競合] occurs when one access (one warp
  instruction) needs more values from a single bank than that bank has
  ports.  The hardware serialises the access over several cycles, one
  per group of values the bank can serve.  In a register file this is a
  \term{register bank conflict}[レジスタバンク競合].
\end{definition}

\Cref{fig:l04-register-bank-conflict} assigns the registers of thread
$k$ to bank $k$ in a four-bank file with two read ports per bank.
\begin{itemize}
  \item Reading \texttt{R0} of threads \texttt{T0}--\texttt{T3}, which
    is what a warp instruction does, touches one entry in each bank:
    four values in one cycle, as fast as the unbanked file.
  \item Reading \texttt{R0}--\texttt{R3} of thread \texttt{T0} touches
    four entries of bank~0, which has two read ports: two cycles.
\end{itemize}

\begin{figure}[htbp]
  \centering
  % Register bank conflicts in a 4-bank file with 2 read ports per bank
% (Lesson 4).  (a) reading R0 of T0..T3 touches one entry in every bank:
% one cycle.  (b) reading R0..R3 of T0 touches four entries of bank 0:
% two cycles.  Redraws the outline figures m4l2_bank_conflict2 and
% m4l2_bank_conflict3.
% Usage: % Register bank conflicts in a 4-bank file with 2 read ports per bank
% (Lesson 4).  (a) reading R0 of T0..T3 touches one entry in every bank:
% one cycle.  (b) reading R0..R3 of T0 touches four entries of bank 0:
% two cycles.  Redraws the outline figures m4l2_bank_conflict2 and
% m4l2_bank_conflict3.
% Usage: % Register bank conflicts in a 4-bank file with 2 read ports per bank
% (Lesson 4).  (a) reading R0 of T0..T3 touches one entry in every bank:
% one cycle.  (b) reading R0..R3 of T0 touches four entries of bank 0:
% two cycles.  Redraws the outline figures m4l2_bank_conflict2 and
% m4l2_bank_conflict3.
% Usage: \input{figures/l04-register-bank-conflict}
% Width: body (120 mm); the drawing is about 118 mm wide.
\begin{tikzpicture}[x=1mm, y=1mm, font=\scriptsize\sffamily]
  % \bankfile{x}{title}{highlight style}: 4 banks x 4 entries
  \newcommand{\bankfile}[2]{%
    \node[anchor=south west, font=\scriptsize\sffamily\bfseries] at (#1,24) {#2};
    \foreach \b in {0,...,3} {
      \node[anchor=south, text=ln-muted] at (#1+\b*12+5,18.5) {\iflangja{バンク \b}{bank \b}};
      \foreach \r in {0,...,3} {
        \draw[fill=ln-box] (#1+\b*12,18-\r*4.2) rectangle ++(10,-4) node[midway] {T\b.R\r};
      }
    }}
  \bankfile{0}{\iflangja{(a) T0--T3 の R0：1 サイクル}{(a) R0 of T0--T3: one cycle}}
  \foreach \b in {0,...,3} {
    \draw[ln-tip, very thick] (\b*12,18) rectangle ++(10,-4);
  }
  \node[anchor=north west, align=left, text=ln-muted] at (0,-0.5)
    {\iflangja{バンクごとに 1 読み出し}{one read per bank}};
  \bankfile{62}{\iflangja{(b) T0 の R0--R3：2 サイクル}{(b) R0--R3 of T0: two cycles}}
  \foreach \r in {0,...,3} {
    \draw[ln-caution, very thick] (62,18-\r*4.2) rectangle ++(10,-4);
  }
  \node[anchor=east, text=ln-caution, align=right] at (61,14.2) {\iflangja{1 周目}{cycle 1}};
  \node[anchor=east, text=ln-caution, align=right] at (61,5.8) {\iflangja{2 周目}{cycle 2}};
  \draw[ln-caution] (61.2,12.1) -- (61.2,18.2);
  \draw[ln-caution] (61.2,3.7) -- (61.2,9.8);
  \node[anchor=north west, align=left, text=ln-muted] at (62,-0.5)
    {\iflangja{バンク 0 に 4 読み出し，ポートは 2 本}{four reads to bank 0, two ports}};
\end{tikzpicture}

% Width: body (120 mm); the drawing is about 118 mm wide.
\begin{tikzpicture}[x=1mm, y=1mm, font=\scriptsize\sffamily]
  % \bankfile{x}{title}{highlight style}: 4 banks x 4 entries
  \newcommand{\bankfile}[2]{%
    \node[anchor=south west, font=\scriptsize\sffamily\bfseries] at (#1,24) {#2};
    \foreach \b in {0,...,3} {
      \node[anchor=south, text=ln-muted] at (#1+\b*12+5,18.5) {\iflangja{バンク \b}{bank \b}};
      \foreach \r in {0,...,3} {
        \draw[fill=ln-box] (#1+\b*12,18-\r*4.2) rectangle ++(10,-4) node[midway] {T\b.R\r};
      }
    }}
  \bankfile{0}{\iflangja{(a) T0--T3 の R0：1 サイクル}{(a) R0 of T0--T3: one cycle}}
  \foreach \b in {0,...,3} {
    \draw[ln-tip, very thick] (\b*12,18) rectangle ++(10,-4);
  }
  \node[anchor=north west, align=left, text=ln-muted] at (0,-0.5)
    {\iflangja{バンクごとに 1 読み出し}{one read per bank}};
  \bankfile{62}{\iflangja{(b) T0 の R0--R3：2 サイクル}{(b) R0--R3 of T0: two cycles}}
  \foreach \r in {0,...,3} {
    \draw[ln-caution, very thick] (62,18-\r*4.2) rectangle ++(10,-4);
  }
  \node[anchor=east, text=ln-caution, align=right] at (61,14.2) {\iflangja{1 周目}{cycle 1}};
  \node[anchor=east, text=ln-caution, align=right] at (61,5.8) {\iflangja{2 周目}{cycle 2}};
  \draw[ln-caution] (61.2,12.1) -- (61.2,18.2);
  \draw[ln-caution] (61.2,3.7) -- (61.2,9.8);
  \node[anchor=north west, align=left, text=ln-muted] at (62,-0.5)
    {\iflangja{バンク 0 に 4 読み出し，ポートは 2 本}{four reads to bank 0, two ports}};
\end{tikzpicture}

% Width: body (120 mm); the drawing is about 118 mm wide.
\begin{tikzpicture}[x=1mm, y=1mm, font=\scriptsize\sffamily]
  % \bankfile{x}{title}{highlight style}: 4 banks x 4 entries
  \newcommand{\bankfile}[2]{%
    \node[anchor=south west, font=\scriptsize\sffamily\bfseries] at (#1,24) {#2};
    \foreach \b in {0,...,3} {
      \node[anchor=south, text=ln-muted] at (#1+\b*12+5,18.5) {\iflangja{バンク \b}{bank \b}};
      \foreach \r in {0,...,3} {
        \draw[fill=ln-box] (#1+\b*12,18-\r*4.2) rectangle ++(10,-4) node[midway] {T\b.R\r};
      }
    }}
  \bankfile{0}{\iflangja{(a) T0--T3 の R0：1 サイクル}{(a) R0 of T0--T3: one cycle}}
  \foreach \b in {0,...,3} {
    \draw[ln-tip, very thick] (\b*12,18) rectangle ++(10,-4);
  }
  \node[anchor=north west, align=left, text=ln-muted] at (0,-0.5)
    {\iflangja{バンクごとに 1 読み出し}{one read per bank}};
  \bankfile{62}{\iflangja{(b) T0 の R0--R3：2 サイクル}{(b) R0--R3 of T0: two cycles}}
  \foreach \r in {0,...,3} {
    \draw[ln-caution, very thick] (62,18-\r*4.2) rectangle ++(10,-4);
  }
  \node[anchor=east, text=ln-caution, align=right] at (61,14.2) {\iflangja{1 周目}{cycle 1}};
  \node[anchor=east, text=ln-caution, align=right] at (61,5.8) {\iflangja{2 周目}{cycle 2}};
  \draw[ln-caution] (61.2,12.1) -- (61.2,18.2);
  \draw[ln-caution] (61.2,3.7) -- (61.2,9.8);
  \node[anchor=north west, align=left, text=ln-muted] at (62,-0.5)
    {\iflangja{バンク 0 に 4 読み出し，ポートは 2 本}{four reads to bank 0, two ports}};
\end{tikzpicture}

  \caption{A four-bank register file with two read ports per bank.
    (a) \texttt{R0} of four threads lies in four banks and is read in
    one cycle.  (b) Four registers of one thread lie in one bank and
    take two cycles.}
  \label{fig:l04-register-bank-conflict}
\end{figure}

The first case looks safe, since a warp always reads the same register
number for all its threads.  It stops being safe once the number of
registers per thread varies, because the thread-to-bank layout then
changes (\cref{fig:l04-register-layout}).  With four registers per
thread and one thread per bank (case~1), the instruction
\texttt{R2 = R0 + R1} reads two values from every bank and completes in
one cycle.  With two registers per thread (case~2) two threads share a
bank: for a warp of \texttt{T0}--\texttt{T3} the same instruction sends
four reads to bank~0 (\texttt{R0} and \texttt{R1} of \texttt{T0} and
\texttt{T1}) and four to bank~1, and takes two cycles while banks~2
and~3 idle.\margin{Real designs spread a thread's registers over the
banks (register $r$ of a thread in bank $r \bmod B$, for instance), so
that the conflicts come from the registers an instruction names, not
from the thread count; the choice of layout is part of the compiler's
and the hardware's contract.}

\begin{figure}[htbp]
  \centering
  % Thread-to-bank layout with a variable number of registers per thread
% (Lesson 4).  Case 1: four registers per thread, one thread per bank.
% Case 2: two registers per thread, two threads per bank; an instruction
% reading R0 and R1 for a warp of T0..T3 sends four reads to each of
% banks 0 and 1.  Redraws the outline figure m4l2_variable (the slide's
% mislabelled third column corrected).
% Usage: % Thread-to-bank layout with a variable number of registers per thread
% (Lesson 4).  Case 1: four registers per thread, one thread per bank.
% Case 2: two registers per thread, two threads per bank; an instruction
% reading R0 and R1 for a warp of T0..T3 sends four reads to each of
% banks 0 and 1.  Redraws the outline figure m4l2_variable (the slide's
% mislabelled third column corrected).
% Usage: % Thread-to-bank layout with a variable number of registers per thread
% (Lesson 4).  Case 1: four registers per thread, one thread per bank.
% Case 2: two registers per thread, two threads per bank; an instruction
% reading R0 and R1 for a warp of T0..T3 sends four reads to each of
% banks 0 and 1.  Redraws the outline figure m4l2_variable (the slide's
% mislabelled third column corrected).
% Usage: \input{figures/l04-register-layout}
% Width: body (120 mm); the drawing is about 118 mm wide.
\begin{tikzpicture}[x=1mm, y=1mm, font=\scriptsize\sffamily]
  \node[anchor=south west, font=\scriptsize\sffamily\bfseries] at (0,26)
    {\iflangja{ケース 1：スレッドあたり 4 レジスタ}{Case 1: four registers per thread}};
  \foreach \b/\c in {0/ln-accent!20, 1/ln-tip!25, 2/ln-caution!20, 3/ln-source!20} {
    \node[anchor=south, text=ln-muted] at (\b*13+5,20.5) {\iflangja{バンク \b}{bank \b}};
    \foreach \r in {0,...,3} {
      \draw[fill=\c] (\b*13,20-\r*4.2) rectangle ++(10,-4) node[midway] {T\b.R\r};
    }
  }
  \node[anchor=north west, align=left, text=ln-muted] at (0,2)
    {\iflangja{R0 と R1 の読み出し：各バンク 2 回}{reading R0 and R1: two reads per bank}};
  \node[anchor=south west, font=\scriptsize\sffamily\bfseries] at (62,26)
    {\iflangja{ケース 2：スレッドあたり 2 レジスタ}{Case 2: two registers per thread}};
  \foreach \b/\c in {0/ln-accent!20, 1/ln-tip!25, 2/ln-caution!20, 3/ln-source!20} {
    \node[anchor=south, text=ln-muted] at (62+\b*13+5,20.5) {\iflangja{バンク \b}{bank \b}};
    \pgfmathtruncatemacro{\ta}{2*\b}
    \pgfmathtruncatemacro{\tb}{2*\b+1}
    \draw[fill=\c] (62+\b*13,20) rectangle ++(10,-4) node[midway] {T\ta.R0};
    \draw[fill=\c] (62+\b*13,15.8) rectangle ++(10,-4) node[midway] {T\ta.R1};
    \draw[fill=\c] (62+\b*13,11.6) rectangle ++(10,-4) node[midway] {T\tb.R0};
    \draw[fill=\c] (62+\b*13,7.4) rectangle ++(10,-4) node[midway] {T\tb.R1};
  }
  \draw[ln-caution, very thick] (61.5,20.5) rectangle (85,2.9);
  \node[anchor=north west, align=left, text=ln-muted] at (62,2)
    {\iflangja{ワープ T0--T3 の R0, R1：\\バンク 0 と 1 に 4 回ずつ}{R0, R1 for warp T0--T3:\\four reads each to banks 0 and 1}};
\end{tikzpicture}

% Width: body (120 mm); the drawing is about 118 mm wide.
\begin{tikzpicture}[x=1mm, y=1mm, font=\scriptsize\sffamily]
  \node[anchor=south west, font=\scriptsize\sffamily\bfseries] at (0,26)
    {\iflangja{ケース 1：スレッドあたり 4 レジスタ}{Case 1: four registers per thread}};
  \foreach \b/\c in {0/ln-accent!20, 1/ln-tip!25, 2/ln-caution!20, 3/ln-source!20} {
    \node[anchor=south, text=ln-muted] at (\b*13+5,20.5) {\iflangja{バンク \b}{bank \b}};
    \foreach \r in {0,...,3} {
      \draw[fill=\c] (\b*13,20-\r*4.2) rectangle ++(10,-4) node[midway] {T\b.R\r};
    }
  }
  \node[anchor=north west, align=left, text=ln-muted] at (0,2)
    {\iflangja{R0 と R1 の読み出し：各バンク 2 回}{reading R0 and R1: two reads per bank}};
  \node[anchor=south west, font=\scriptsize\sffamily\bfseries] at (62,26)
    {\iflangja{ケース 2：スレッドあたり 2 レジスタ}{Case 2: two registers per thread}};
  \foreach \b/\c in {0/ln-accent!20, 1/ln-tip!25, 2/ln-caution!20, 3/ln-source!20} {
    \node[anchor=south, text=ln-muted] at (62+\b*13+5,20.5) {\iflangja{バンク \b}{bank \b}};
    \pgfmathtruncatemacro{\ta}{2*\b}
    \pgfmathtruncatemacro{\tb}{2*\b+1}
    \draw[fill=\c] (62+\b*13,20) rectangle ++(10,-4) node[midway] {T\ta.R0};
    \draw[fill=\c] (62+\b*13,15.8) rectangle ++(10,-4) node[midway] {T\ta.R1};
    \draw[fill=\c] (62+\b*13,11.6) rectangle ++(10,-4) node[midway] {T\tb.R0};
    \draw[fill=\c] (62+\b*13,7.4) rectangle ++(10,-4) node[midway] {T\tb.R1};
  }
  \draw[ln-caution, very thick] (61.5,20.5) rectangle (85,2.9);
  \node[anchor=north west, align=left, text=ln-muted] at (62,2)
    {\iflangja{ワープ T0--T3 の R0, R1：\\バンク 0 と 1 に 4 回ずつ}{R0, R1 for warp T0--T3:\\four reads each to banks 0 and 1}};
\end{tikzpicture}

% Width: body (120 mm); the drawing is about 118 mm wide.
\begin{tikzpicture}[x=1mm, y=1mm, font=\scriptsize\sffamily]
  \node[anchor=south west, font=\scriptsize\sffamily\bfseries] at (0,26)
    {\iflangja{ケース 1：スレッドあたり 4 レジスタ}{Case 1: four registers per thread}};
  \foreach \b/\c in {0/ln-accent!20, 1/ln-tip!25, 2/ln-caution!20, 3/ln-source!20} {
    \node[anchor=south, text=ln-muted] at (\b*13+5,20.5) {\iflangja{バンク \b}{bank \b}};
    \foreach \r in {0,...,3} {
      \draw[fill=\c] (\b*13,20-\r*4.2) rectangle ++(10,-4) node[midway] {T\b.R\r};
    }
  }
  \node[anchor=north west, align=left, text=ln-muted] at (0,2)
    {\iflangja{R0 と R1 の読み出し：各バンク 2 回}{reading R0 and R1: two reads per bank}};
  \node[anchor=south west, font=\scriptsize\sffamily\bfseries] at (62,26)
    {\iflangja{ケース 2：スレッドあたり 2 レジスタ}{Case 2: two registers per thread}};
  \foreach \b/\c in {0/ln-accent!20, 1/ln-tip!25, 2/ln-caution!20, 3/ln-source!20} {
    \node[anchor=south, text=ln-muted] at (62+\b*13+5,20.5) {\iflangja{バンク \b}{bank \b}};
    \pgfmathtruncatemacro{\ta}{2*\b}
    \pgfmathtruncatemacro{\tb}{2*\b+1}
    \draw[fill=\c] (62+\b*13,20) rectangle ++(10,-4) node[midway] {T\ta.R0};
    \draw[fill=\c] (62+\b*13,15.8) rectangle ++(10,-4) node[midway] {T\ta.R1};
    \draw[fill=\c] (62+\b*13,11.6) rectangle ++(10,-4) node[midway] {T\tb.R0};
    \draw[fill=\c] (62+\b*13,7.4) rectangle ++(10,-4) node[midway] {T\tb.R1};
  }
  \draw[ln-caution, very thick] (61.5,20.5) rectangle (85,2.9);
  \node[anchor=north west, align=left, text=ln-muted] at (62,2)
    {\iflangja{ワープ T0--T3 の R0, R1：\\バンク 0 と 1 に 4 回ずつ}{R0, R1 for warp T0--T3:\\four reads each to banks 0 and 1}};
\end{tikzpicture}

  \caption{Thread-to-bank layout with a variable register count.
    Case~1: four registers per thread, one thread per bank.  Case~2:
    two registers per thread, two threads per bank, so an instruction
    reading \texttt{R0} and \texttt{R1} for warp \texttt{T0}--\texttt{T3}
    sends four reads to each of banks~0 and~1.}
  \label{fig:l04-register-layout}
\end{figure}

\subsection{Static and dynamic}

The first remedy is the compiler.  Register numbers are fixed in the
program text, so the compiler can choose which registers an instruction
uses and in what order instructions are emitted so that the operands of
one instruction fall into different banks.  It cannot remove every
conflict, because the dependences of the program constrain the order and
the number of registers constrains the choice.\margin{Register
allocation, the compiler pass that assigns program values to registers,
builds on the liveness analysis of Lesson~11; on a GPU the bank layout
is one of its inputs.}  The distinction behind this remedy is used throughout the
compiler lessons.

\begin{definition}[Static and dynamic]
  \label{def:l04-static-dynamic}
  A property of a program is \term{static}[静的] when it is determined
  before the program runs and does not depend on its input; it is
  \term{dynamic}[動的] when it depends on the input and is known only
  during execution.  Static analysis is compile-time analysis.  The
  \term{static instruction count}[静的命令数] is the number of
  instructions in the program text; the
  \term{dynamic instruction count}[動的命令数] is the number of
  instructions executed in one run.
\end{definition}

\begin{example}
  \label{ex:l04-static-dynamic}
  The loop below sums 16 words and consists of five instructions.
\begin{lstlisting}[numbers=none]
LOOP: LD   R2, [R3]      // load a word
      ADD  R1, R1, R2    // accumulate
      ADD  R3, R3, #4    // next address
      SUB  R4, R4, #1    // count down from 16
      BNEZ R4, LOOP      // repeat while R4 != 0
\end{lstlisting}
  Its static instruction count is 5.  With \texttt{R4} initialised to 16
  the loop body runs 16 times, so the dynamic instruction count is
  $16 \times 5 = 80$; with \texttt{R4} set from the input to 1000 it is
  \num{5000}.  The register numbers, and therefore the banks each
  instruction reads, are static; the trip count is dynamic.
\end{example}

\subsection{Operand collectors and the scoreboard}

The second remedy is in the pipeline.  Because a register read may take
several cycles, through a bank conflict or because a bank has a single
read port, the SM does not read operands and execute in consecutive
stages as the five-stage CPU does.  \Cref{fig:l04-pipeline} inserts two
things between register read and execution.

\begin{definition}[Operand buffer and operand collector]
  \label{def:l04-operand-collector}
  An \term{operand buffer}[オペランドバッファ] is storage, one entry per
  waiting warp instruction and per lane, into which source operands are
  written as the register banks deliver them over several cycles.  The
  \term{operand collector}[オペランドコレクタ] is the unit that owns these
  buffers, issues the bank reads for the instructions it holds, and
  reports an instruction as ready when all of its operands have
  arrived.\margin{\emph{Operand collector} is the term of A18 ch.~3
  \cite{aamodt2018gpu}; K12 §1.4.1 describes the same structure as
  operand buffers and notes that it also serves as a result queue
  \cite{kim2012gpgpu}.}
\end{definition}

\begin{figure}[htbp]
  \centering
  % The SM pipeline with register banks, an operand collector and a
% scoreboard (Lesson 4): fetch (one PC per warp, I-cache), decode,
% register-file read from banks, the operand buffer / collector, the
% scheduler with its scoreboard, execution (SIMT lanes, SFUs, memory and
% shared memory) and write-back.  Redraws the outline figures
% m4l2_solution and m4l3_gpu_pipeline.
% Usage: % The SM pipeline with register banks, an operand collector and a
% scoreboard (Lesson 4): fetch (one PC per warp, I-cache), decode,
% register-file read from banks, the operand buffer / collector, the
% scheduler with its scoreboard, execution (SIMT lanes, SFUs, memory and
% shared memory) and write-back.  Redraws the outline figures
% m4l2_solution and m4l3_gpu_pipeline.
% Usage: % The SM pipeline with register banks, an operand collector and a
% scoreboard (Lesson 4): fetch (one PC per warp, I-cache), decode,
% register-file read from banks, the operand buffer / collector, the
% scheduler with its scoreboard, execution (SIMT lanes, SFUs, memory and
% shared memory) and write-back.  Redraws the outline figures
% m4l2_solution and m4l3_gpu_pipeline.
% Usage: \input{figures/l04-pipeline}
% Width: body (120 mm); the drawing is about 118 mm wide.
\begin{tikzpicture}[x=1mm, y=1mm, font=\scriptsize\sffamily,
    >={Stealth[length=1.6mm]},
    st/.style={draw, fill=ln-box, minimum height=8mm, minimum width=15mm,
               align=center, inner sep=2pt},
    mem/.style={draw, fill=ln-accent!15, minimum height=6mm, align=center,
                inner sep=2pt}]
  \node[st] (fetch) at (7.5,20) {\iflangja{フェッチ}{Fetch}};
  \node[st, right=2mm of fetch, minimum width=13mm] (dec) {\iflangja{デコード}{Decode}};
  \node[st, right=2mm of dec, minimum width=14mm] (reg)
    {\iflangja{レジスタ\\読み出し}{Register\\read}};
  \node[st, right=2mm of reg, minimum width=16mm] (oc)
    {\iflangja{オペランド\\コレクタ}{Operand\\collector}};
  \node[st, right=2mm of oc, minimum width=15mm] (sched)
    {\iflangja{スケジュール}{Schedule}};
  \node[st, right=2mm of sched, minimum width=14mm, fill=ln-accent!12] (ex)
    {\iflangja{実行}{Execute}};
  \draw[->] (fetch) -- (dec);
  \draw[->] (dec) -- (reg);
  \draw[->] (reg) -- (oc);
  \draw[->] (oc) -- (sched);
  \draw[->] (sched) -- (ex);
  % write-back below execute, result back to the register banks
  \node[st, below=9mm of ex, minimum width=14mm] (wb)
    {\iflangja{ライトバック}{Write-back}};
  \draw[->] (ex) -- (wb);
  % state attached to the stages
  \node[mem, above=4mm of fetch, minimum width=15mm] (pcs)
    {\iflangja{ワープごとの PC}{one PC per warp}};
  \draw[<->] (pcs) -- (fetch);
  \node[mem, below=4mm of fetch, minimum width=15mm] (icache)
    {\iflangja{命令キャッシュ}{I-cache}};
  \draw[<->] (icache) -- (fetch);
  \node[mem, below=4mm of reg, minimum width=14mm] (banks)
    {\iflangja{レジスタバンク\\0 1 2 3}{register banks\\0 1 2 3}};
  \draw[<->] (banks) -- (reg);
  \draw[->] (wb.west) -| (banks.south);
  \node[mem, above=4mm of oc, minimum width=16mm] (buf)
    {\iflangja{オペランド\\バッファ}{operand\\buffer}};
  \draw[<->] (buf) -- (oc);
  \node[mem, above=4mm of sched, minimum width=15mm] (sb)
    {\iflangja{スコアボード}{scoreboard}};
  \draw[<->] (sb) -- (sched);
  \node[mem, above=4mm of ex, minimum width=14mm] (lanes)
    {\iflangja{レーン，SFU}{lanes, SFUs}};
  \draw[<->] (lanes) -- (ex);
  \node[mem, right=2mm of ex, minimum width=11mm] (dmem)
    {\iflangja{メモリ\\共有メモリ}{memory,\\shared mem.}};
  \draw[<->] (ex) -- (dmem);
\end{tikzpicture}

% Width: body (120 mm); the drawing is about 118 mm wide.
\begin{tikzpicture}[x=1mm, y=1mm, font=\scriptsize\sffamily,
    >={Stealth[length=1.6mm]},
    st/.style={draw, fill=ln-box, minimum height=8mm, minimum width=15mm,
               align=center, inner sep=2pt},
    mem/.style={draw, fill=ln-accent!15, minimum height=6mm, align=center,
                inner sep=2pt}]
  \node[st] (fetch) at (7.5,20) {\iflangja{フェッチ}{Fetch}};
  \node[st, right=2mm of fetch, minimum width=13mm] (dec) {\iflangja{デコード}{Decode}};
  \node[st, right=2mm of dec, minimum width=14mm] (reg)
    {\iflangja{レジスタ\\読み出し}{Register\\read}};
  \node[st, right=2mm of reg, minimum width=16mm] (oc)
    {\iflangja{オペランド\\コレクタ}{Operand\\collector}};
  \node[st, right=2mm of oc, minimum width=15mm] (sched)
    {\iflangja{スケジュール}{Schedule}};
  \node[st, right=2mm of sched, minimum width=14mm, fill=ln-accent!12] (ex)
    {\iflangja{実行}{Execute}};
  \draw[->] (fetch) -- (dec);
  \draw[->] (dec) -- (reg);
  \draw[->] (reg) -- (oc);
  \draw[->] (oc) -- (sched);
  \draw[->] (sched) -- (ex);
  % write-back below execute, result back to the register banks
  \node[st, below=9mm of ex, minimum width=14mm] (wb)
    {\iflangja{ライトバック}{Write-back}};
  \draw[->] (ex) -- (wb);
  % state attached to the stages
  \node[mem, above=4mm of fetch, minimum width=15mm] (pcs)
    {\iflangja{ワープごとの PC}{one PC per warp}};
  \draw[<->] (pcs) -- (fetch);
  \node[mem, below=4mm of fetch, minimum width=15mm] (icache)
    {\iflangja{命令キャッシュ}{I-cache}};
  \draw[<->] (icache) -- (fetch);
  \node[mem, below=4mm of reg, minimum width=14mm] (banks)
    {\iflangja{レジスタバンク\\0 1 2 3}{register banks\\0 1 2 3}};
  \draw[<->] (banks) -- (reg);
  \draw[->] (wb.west) -| (banks.south);
  \node[mem, above=4mm of oc, minimum width=16mm] (buf)
    {\iflangja{オペランド\\バッファ}{operand\\buffer}};
  \draw[<->] (buf) -- (oc);
  \node[mem, above=4mm of sched, minimum width=15mm] (sb)
    {\iflangja{スコアボード}{scoreboard}};
  \draw[<->] (sb) -- (sched);
  \node[mem, above=4mm of ex, minimum width=14mm] (lanes)
    {\iflangja{レーン，SFU}{lanes, SFUs}};
  \draw[<->] (lanes) -- (ex);
  \node[mem, right=2mm of ex, minimum width=11mm] (dmem)
    {\iflangja{メモリ\\共有メモリ}{memory,\\shared mem.}};
  \draw[<->] (ex) -- (dmem);
\end{tikzpicture}

% Width: body (120 mm); the drawing is about 118 mm wide.
\begin{tikzpicture}[x=1mm, y=1mm, font=\scriptsize\sffamily,
    >={Stealth[length=1.6mm]},
    st/.style={draw, fill=ln-box, minimum height=8mm, minimum width=15mm,
               align=center, inner sep=2pt},
    mem/.style={draw, fill=ln-accent!15, minimum height=6mm, align=center,
                inner sep=2pt}]
  \node[st] (fetch) at (7.5,20) {\iflangja{フェッチ}{Fetch}};
  \node[st, right=2mm of fetch, minimum width=13mm] (dec) {\iflangja{デコード}{Decode}};
  \node[st, right=2mm of dec, minimum width=14mm] (reg)
    {\iflangja{レジスタ\\読み出し}{Register\\read}};
  \node[st, right=2mm of reg, minimum width=16mm] (oc)
    {\iflangja{オペランド\\コレクタ}{Operand\\collector}};
  \node[st, right=2mm of oc, minimum width=15mm] (sched)
    {\iflangja{スケジュール}{Schedule}};
  \node[st, right=2mm of sched, minimum width=14mm, fill=ln-accent!12] (ex)
    {\iflangja{実行}{Execute}};
  \draw[->] (fetch) -- (dec);
  \draw[->] (dec) -- (reg);
  \draw[->] (reg) -- (oc);
  \draw[->] (oc) -- (sched);
  \draw[->] (sched) -- (ex);
  % write-back below execute, result back to the register banks
  \node[st, below=9mm of ex, minimum width=14mm] (wb)
    {\iflangja{ライトバック}{Write-back}};
  \draw[->] (ex) -- (wb);
  % state attached to the stages
  \node[mem, above=4mm of fetch, minimum width=15mm] (pcs)
    {\iflangja{ワープごとの PC}{one PC per warp}};
  \draw[<->] (pcs) -- (fetch);
  \node[mem, below=4mm of fetch, minimum width=15mm] (icache)
    {\iflangja{命令キャッシュ}{I-cache}};
  \draw[<->] (icache) -- (fetch);
  \node[mem, below=4mm of reg, minimum width=14mm] (banks)
    {\iflangja{レジスタバンク\\0 1 2 3}{register banks\\0 1 2 3}};
  \draw[<->] (banks) -- (reg);
  \draw[->] (wb.west) -| (banks.south);
  \node[mem, above=4mm of oc, minimum width=16mm] (buf)
    {\iflangja{オペランド\\バッファ}{operand\\buffer}};
  \draw[<->] (buf) -- (oc);
  \node[mem, above=4mm of sched, minimum width=15mm] (sb)
    {\iflangja{スコアボード}{scoreboard}};
  \draw[<->] (sb) -- (sched);
  \node[mem, above=4mm of ex, minimum width=14mm] (lanes)
    {\iflangja{レーン，SFU}{lanes, SFUs}};
  \draw[<->] (lanes) -- (ex);
  \node[mem, right=2mm of ex, minimum width=11mm] (dmem)
    {\iflangja{メモリ\\共有メモリ}{memory,\\shared mem.}};
  \draw[<->] (ex) -- (dmem);
\end{tikzpicture}

  \caption{The SM pipeline of this lesson: fetch with one PC per warp,
    decode, register read from the banks, the operand collector with
    its buffers, the scheduler with its scoreboard, execution on the
    lanes and SFUs or in the memory units, and write-back to the
    banks.}
  \label{fig:l04-pipeline}
\end{figure}

\begin{definition}[Scoreboard]
  \label{def:l04-scoreboard}
  A \term{scoreboard}[スコアボード] is a table that records, for each
  pending warp instruction, whether each of its source operands is
  ready.  The scheduler consults it to select a warp whose operands are
  all ready and issues that warp's instruction to the lanes; among
  several ready warps it applies a policy such as oldest first.\caution{Lecture
  vs textbook: in A18 ch.~3 and K12 §1.4.1 the scoreboard records which
  registers have pending writes and blocks a warp's instruction on a RAW
  or WAW hazard before issue; operand collection follows issue, so the
  hardware order is decode, scoreboard check and issue, operand
  collection, execute \cite{aamodt2018gpu,kim2012gpgpu}.}
\end{definition}

On a CPU the scoreboard was the mechanism of early out-of-order dynamic
scheduling: it tracked the hazards among the instructions of one thread.
On a GPU it answers a simpler question, whether a warp instruction has
its operands, and the choice is among warps rather than among the
instructions of one thread.

\begin{example}
  \label{ex:l04-scoreboard}
  Four warps \texttt{w0}--\texttt{w3} each have a two-operand
  instruction waiting in the operand collector, and the banks deliver
  one operand per cycle in the order shown.  The scoreboard columns give
  the readiness of \texttt{Src1} and \texttt{Src2} after each cycle.
  \begin{center}
    \begin{tabular}{clcccc}
      \toprule
      cycle & operand delivered & \texttt{w0} & \texttt{w1} & \texttt{w2} & \texttt{w3} \\
      \midrule
      0 & --- & $\cdot\,\cdot$ & $\cdot\,\cdot$ & $\cdot\,\cdot$ & $\cdot\,\cdot$ \\
      1 & \texttt{w0.Src1} & $\bullet\,\cdot$ & $\cdot\,\cdot$ & $\cdot\,\cdot$ & $\cdot\,\cdot$ \\
      2 & \texttt{w1.Src2} & $\bullet\,\cdot$ & $\cdot\,\bullet$ & $\cdot\,\cdot$ & $\cdot\,\cdot$ \\
      3 & \texttt{w2.Src1} & $\bullet\,\cdot$ & $\cdot\,\bullet$ & $\bullet\,\cdot$ & $\cdot\,\cdot$ \\
      4 & \texttt{w0.Src2} & $\bullet\,\bullet$ & $\cdot\,\bullet$ & $\bullet\,\cdot$ & $\cdot\,\cdot$ \\
      \bottomrule
    \end{tabular}
  \end{center}
  After cycle~4 \texttt{w0} is the only warp with both operands ready,
  so the scheduler selects it and its operands move from the buffers
  into the lanes; \texttt{w1} and \texttt{w2} keep waiting for their
  second operand.  The order in which operands arrive, and hence which
  warp issues first, is dynamic: it depends on bank conflicts and on
  the memory system, not on the program text.\sidenote{The scoreboard was introduced with
the CDC~6600 (1964) for a single instruction stream.  The GPU form is
in-order per warp: K12 §1.4.1 observes that the scheduler is in order
but \enquote{can select any warps that are ready}, so that the
execution of the SM as a whole \enquote{might look like an out-of-order
execution} \cite{kim2012gpgpu}.  Lesson~7 compares the selection
policies (round-robin, greedy-then-oldest) in detail.}
\end{example}

\subsection{Shared-memory bank conflicts}

Shared memory, the per-block scratchpad of Lesson~3, is banked for the
same reason as the register file: a warp's load must deliver 32 words
in a cycle.  \Cref{fig:l04-shared-banks} shows a four-bank shared memory
whose consecutive words are interleaved across the banks, word $a$ in
bank $a \bmod 4$.  If thread $k$ reads
\texttt{sharedInput[threadIdx.x * 4]}, that is word $4k$, every thread
of the warp reads bank~0 and the load is serialised over as many cycles
as there are threads.

\begin{marginfigure}
  % A four-bank shared memory (Lesson 4): word addresses interleaved over
% the banks (bank = address mod 4), and four threads reading
% sharedInput[threadIdx.x * 4], which all hit bank 0.  Redraws the
% outline figure m4l2_shared_mem.
% Usage: % A four-bank shared memory (Lesson 4): word addresses interleaved over
% the banks (bank = address mod 4), and four threads reading
% sharedInput[threadIdx.x * 4], which all hit bank 0.  Redraws the
% outline figure m4l2_shared_mem.
% Usage: % A four-bank shared memory (Lesson 4): word addresses interleaved over
% the banks (bank = address mod 4), and four threads reading
% sharedInput[threadIdx.x * 4], which all hit bank 0.  Redraws the
% outline figure m4l2_shared_mem.
% Usage: \input{figures/l04-shared-banks} inside marginfigure
% Width: margin (50 mm); the drawing is about 48 mm wide.
\begin{tikzpicture}[x=1mm, y=1mm, font=\scriptsize\sffamily,
    >={Stealth[length=1.4mm]}]
  \foreach \b in {0,...,3} {
    \node[anchor=south, text=ln-muted] at (14+\b*9+3.5,21) {\iflangja{\b}{bank \b}};
    \foreach \k in {0,...,3} {
      \pgfmathtruncatemacro{\a}{\k*4+\b}
      \draw[fill=ln-box] (14+\b*9,20-\k*4.5) rectangle ++(7,-4) node[midway] {\a};
    }
  }
  \iflangja{\node[anchor=south, text=ln-muted] at (31,24) {バンク};}{}% JA: bank numbers only, one heading
  \foreach \k in {0,...,3} {
    \node[anchor=east, inner sep=1pt] (t\k) at (8,18-\k*4.5) {T\k};
    \draw[->, ln-caution, thick] (t\k.east) -- (14,18-\k*4.5);
  }
  \node[anchor=north west, align=left, text=ln-muted] at (0,0.5)
    {\iflangja{T$k$ は語 $4k$ を読む：全員がバンク 0}{T$k$ reads word $4k$: all in bank 0}};
\end{tikzpicture}
 inside marginfigure
% Width: margin (50 mm); the drawing is about 48 mm wide.
\begin{tikzpicture}[x=1mm, y=1mm, font=\scriptsize\sffamily,
    >={Stealth[length=1.4mm]}]
  \foreach \b in {0,...,3} {
    \node[anchor=south, text=ln-muted] at (14+\b*9+3.5,21) {\iflangja{\b}{bank \b}};
    \foreach \k in {0,...,3} {
      \pgfmathtruncatemacro{\a}{\k*4+\b}
      \draw[fill=ln-box] (14+\b*9,20-\k*4.5) rectangle ++(7,-4) node[midway] {\a};
    }
  }
  \iflangja{\node[anchor=south, text=ln-muted] at (31,24) {バンク};}{}% JA: bank numbers only, one heading
  \foreach \k in {0,...,3} {
    \node[anchor=east, inner sep=1pt] (t\k) at (8,18-\k*4.5) {T\k};
    \draw[->, ln-caution, thick] (t\k.east) -- (14,18-\k*4.5);
  }
  \node[anchor=north west, align=left, text=ln-muted] at (0,0.5)
    {\iflangja{T$k$ は語 $4k$ を読む：全員がバンク 0}{T$k$ reads word $4k$: all in bank 0}};
\end{tikzpicture}
 inside marginfigure
% Width: margin (50 mm); the drawing is about 48 mm wide.
\begin{tikzpicture}[x=1mm, y=1mm, font=\scriptsize\sffamily,
    >={Stealth[length=1.4mm]}]
  \foreach \b in {0,...,3} {
    \node[anchor=south, text=ln-muted] at (14+\b*9+3.5,21) {\iflangja{\b}{bank \b}};
    \foreach \k in {0,...,3} {
      \pgfmathtruncatemacro{\a}{\k*4+\b}
      \draw[fill=ln-box] (14+\b*9,20-\k*4.5) rectangle ++(7,-4) node[midway] {\a};
    }
  }
  \iflangja{\node[anchor=south, text=ln-muted] at (31,24) {バンク};}{}% JA: bank numbers only, one heading
  \foreach \k in {0,...,3} {
    \node[anchor=east, inner sep=1pt] (t\k) at (8,18-\k*4.5) {T\k};
    \draw[->, ln-caution, thick] (t\k.east) -- (14,18-\k*4.5);
  }
  \node[anchor=north west, align=left, text=ln-muted] at (0,0.5)
    {\iflangja{T$k$ は語 $4k$ を読む：全員がバンク 0}{T$k$ reads word $4k$: all in bank 0}};
\end{tikzpicture}

  \caption{Four shared-memory banks with interleaved words; a stride of
    four words sends every thread to bank~0.}
  \label{fig:l04-shared-banks}
\end{marginfigure}

\begin{definition}[Shared-memory bank conflict]
  \label{def:l04-smem-bank-conflict}
  A \term{shared memory bank conflict}[共有メモリのバンク競合] occurs when
  the threads of one warp access different words that lie in the same
  shared-memory bank in one instruction.  The hardware splits the
  request into as many conflict-free requests as necessary, and the
  effective bandwidth falls by that factor.  Threads that access the
  \emph{same} word do not conflict: a read of one word by several
  threads is served by a broadcast, and a write of one word by several
  threads is performed by one of them (BPG §10.2.3.1 Shared Memory and
  Memory Banks \cite{nvidia2026cudabestpractices}; CPG, Compute
  Capability 5.x, Shared Memory \cite{nvidia2024cudaguide}).
\end{definition}

On NVIDIA GPUs of compute capability 5.0 and later, shared memory has
32 banks, each \SI{32}{\bit} wide and delivering one word per clock,
and successive 32-bit words are assigned to successive banks.  With a
warp of 32 threads a conflict can therefore arise between any two
threads of the warp.\margin{Shared-memory latency is \enquote{roughly
100x lower than uncached global memory latency}
\cite{harris2013sharedmem}; an $n$-way conflict divides the effective
bandwidth by $n$ (BPG §10.2.3.1).}  From these rules the bank of a byte address
$a$ follows as
\begin{equation}
  \text{bank}(a) = \left\lfloor a / 4 \right\rfloor \bmod 32,
  \label{eq:l04-bank}
\end{equation}
and for 32-bit elements accessed at a constant
\term{stride}[ストライド] of $s$ words, the distance between the
elements read by consecutive threads (thread $k$ reads element $sk$),
the banks touched
repeat after $32 / \gcd(s, 32)$ threads, so the access is a
$\gcd(s, 32)$-way conflict.

\begin{example}
  \label{ex:l04-smem-stride}
  A warp reads \texttt{s[threadIdx.x * 8]} from an array of
  \texttt{float}.  The word stride is $s = 8$, so thread $k$ reads bank
  $8k \bmod 32$: threads 0, 4, 8, \ldots, 28 all hit bank~0, threads
  1, 5, \ldots, 29 bank~8, and only banks 0, 8, 16 and 24 are used.
  Eight threads share each bank, so the load is an 8-way conflict,
  $\gcd(8, 32) = 8$, and takes eight shared-memory cycles instead of
  one.  With stride 7 the banks are $7k \bmod 32$, all 32 distinct, and
  the load takes one cycle.\caution{Lecture vs guide: the four-bank
  model turns \texttt{threadIdx.x * 4} into a full conflict (every
  thread in one bank).  On real 32-bank hardware the same access is a
  4-way conflict, $\gcd(4, 32) = 4$, using eight banks.}
\end{example}

The common case in practice is a tile stored row by row and read column
by column, as in the tiled matrix multiplication of Lesson~5.  With a
$32 \times 32$ tile of 32-bit elements, the elements of one column are
32 words apart, a 32-way conflict.  The standard fix is
\term{padding}[パディング]: declaring the tile as $32 \times 33$ makes
the column stride 33, which by the modulo arithmetic of
\cref{eq:l04-bank} behaves like a unit stride and is conflict-free, at
the cost of one unused word per row.\margin{BPG §10.2.3.3 Shared
Memory in Matrix Multiplication (C = AA$^\mathsf{T}$) is the source of
the $[w][w+1]$ padding rule \cite{nvidia2026cudabestpractices}; Harris
discusses the same conflict for transposes \cite{harris2013sharedmem}.}
\begin{lstlisting}[language=C++, numbers=none]
__shared__ float tileA[32][32];  // column read: 32-way
__shared__ float tileB[32][33];  // column read: no conflict
// (tx, ty) reads tileA[tx][ty] / tileB[tx][ty]:
// stride 32 vs 33
\end{lstlisting}
Restructuring the data layout or the access pattern in this way is the
software remedy for shared-memory conflicts, and the compiler cannot
apply it automatically because the addresses are dynamic.\tip{Any odd word stride is
conflict-free on 32 banks, since $\gcd(s, 32) = 1$; padding a row by
one word makes an even row stride odd.}

\begin{keypoint}
  Banking gives a wide register file or shared memory the bandwidth of
  many ports at the cost of a few: values in different banks are read
  together, values in the same bank are serialised.  Register bank
  conflicts are absorbed by the operand collector and the scoreboard,
  which lets the scheduler issue whichever warp has its operands;
  shared-memory conflicts are removed by the programmer through strides
  and padding.
\end{keypoint}

\section{GPU architecture pipeline}
\label{sec:l04-pipeline}

This section follows warp instructions through the pipeline of
\cref{fig:l04-pipeline} on a small configuration, stage by stage, and
then adds the mask bits that let a warp instruction skip some of its
threads.\source{M4, pipeline walk-through; K12 §1.4.1; A18 ch.~3; PTX
ISA, special registers \cite{nvidia2026ptx}.}

\subsection{The stages}

The stages of \cref{fig:l04-pipeline} are, for one warp instruction:
\begin{description}
  \item[Fetch.]  The \term{instruction fetch}[命令フェッチ] stage reads the
    instruction at the selected warp's PC from the instruction cache,
    one instruction for the whole warp.
  \item[Decode.]  The \term{decode}[デコード] stage identifies the operation, the
    source and destination registers and any immediate operand; an
    immediate is the same for every thread and is broadcast into every
    lane's operand buffer at once, already marked ready.
  \item[Register read.]  The register banks deliver the source values of
    each thread into the operand buffers, over one or more cycles.
  \item[Schedule and issue.]  When the scoreboard shows all operands
    ready the scheduler selects the warp; \term{issue}[命令発行] sends
    the instruction and its operands to the execution units.
  \item[Execute.]  The arithmetic lanes compute, one thread per lane.
    Transcendental functions (sine, reciprocal, square root) go to a
    \term{special function unit (SFU)}[特殊関数ユニット]\index{SFU|see{special function unit (SFU)}}, of which an SM
    has a few, so that a warp's SFU instruction takes several passes;
    loads and stores go to the memory unit or to shared
    memory.\sidenote{K12 §1.4.1 gives the shape of the execution stage
    on G80: 8 lanes, so a 32-thread warp took four cycles per
    instruction; 2--4 SFUs per SM computing approximations of the
    transcendental functions; and no forwarding network as in a CPU
    pipeline, since a warp result is as wide as the whole warp (K12
    puts it at \SI{1}{\kibi\byte}), so that dependent back-to-back
    instructions rely on the other warps to fill the gap
    \cite{kim2012gpgpu}.}
  \item[Write-back.]  The \term{write-back}[ライトバック] stage writes each
    thread's result into its register in the banks.
\end{description}

\subsection{A walk-through}

The configuration of the example is deliberately small: a warp size of
4, two thread blocks of 8 threads, hence four warps, and an SM with four
lanes.  Blocks and threads are numbered from 0 as in CUDA, so
\texttt{\%ctaid.x} is 0 or 1 and \texttt{\%tid.x} runs from 0 to 7;
registers are 32-bit and their initial values are arbitrary.  The state
before the first cycle is given by the three tables of
\cref{tab:l04-walkthrough-state}.

\begin{table}[htbp]
  \centering
  \caption{Initial state of the walk-through: the PC of each warp, the
    instructions in the instruction cache, and the register file of
    block~0.}
  \label{tab:l04-walkthrough-state}
  \begin{tabular}{lll}
    \toprule
    warp & block, threads & PC \\
    \midrule
    \texttt{W0} & block 0, \texttt{T0}--\texttt{T3} & \texttt{0x8000} \\
    \texttt{W1} & block 0, \texttt{T4}--\texttt{T7} & \texttt{0x8004} \\
    \texttt{W2} & block 1, \texttt{T0}--\texttt{T3} & \texttt{0x8000} \\
    \texttt{W3} & block 1, \texttt{T4}--\texttt{T7} & \texttt{0x800C} \\
    \bottomrule
  \end{tabular}
  \hspace{6mm}
  \begin{tabular}{ll}
    \toprule
    address & instruction \\
    \midrule
    \texttt{0x8000} & \texttt{add r1, r1, 2} \\
    \texttt{0x8004} & \texttt{sub r2, r2, 1} \\
    \texttt{0x8008} & \texttt{mov r1, \%tid.x} \\
    \texttt{0x800C} & \texttt{mov r2, \%ctaid.x} \\
    \texttt{0x8010} & \texttt{ld.shared r1, [A]} \\
    \bottomrule
  \end{tabular}

  \medskip
  \begin{tabular}{lcccccccc}
    \toprule
    block 0 & \texttt{T0} & \texttt{T1} & \texttt{T2} & \texttt{T3} & \texttt{T4} & \texttt{T5} & \texttt{T6} & \texttt{T7} \\
    \midrule
    \texttt{R1} & 20 & 21 & 22 & 23 & 30 & 31 & 32 & 33 \\
    \texttt{R2} &  5 &  5 &  5 &  5 &  7 &  7 &  7 &  7 \\
    \texttt{R3} &  1 &  4 &  9 & 16 & 25 & 36 & 49 & 64 \\
    \texttt{R4} &  8 &  6 &  7 &  5 &  3 &  0 &  9 &  2 \\
    \bottomrule
  \end{tabular}
\end{table}

The instructions use two \term{special register}[特殊レジスタ]s of PTX,
read-only registers that the hardware fills for every thread:
\texttt{\%tid.x} holds the thread's index within its block and
\texttt{\%ctaid.x} the index of its block (CTA) within the
grid.\margin{PTX ISA 9.4, special registers \texttt{\%tid},
\texttt{\%ntid}, \texttt{\%ctaid}, \texttt{\%nctaid}
\cite{nvidia2026ptx}; they are what \texttt{threadIdx}, \texttt{blockDim}
and \texttt{blockIdx} of Lesson~3 compile to.  Lesson~10 presents PTX
as a whole.}  The scheduler is assumed to pick warps in the order
\texttt{W0}, \texttt{W1}, \texttt{W3}, and each step below takes the
instruction through all stages before the next warp is considered.
\begin{enumerate}
  \item \emph{\texttt{W0} at \texttt{0x8000}: \texttt{add r1, r1, 2}.}
    Fetch reads one instruction for the four threads.  Decode broadcasts
    the immediate 2 into \texttt{Src2} of all four lanes and marks it
    ready.  Register read fills \texttt{Src1} with \texttt{R1} of
    \texttt{T0}--\texttt{T3}, that is 20, 21, 22, 23.  Both operands
    are ready, the scoreboard lets the scheduler issue \texttt{W0}, the
    four lanes add, and write-back stores 22, 23, 24, 25 into
    \texttt{R1} of \texttt{T0}--\texttt{T3}.  The PC of \texttt{W0}
    advances to \texttt{0x8004}.
  \item \emph{\texttt{W1} at \texttt{0x8004}: \texttt{sub r2, r2, 1}.}
    The same stages for threads \texttt{T4}--\texttt{T7}: the immediate
    1 is broadcast, \texttt{R2} of \texttt{T4}--\texttt{T7} is read as
    7, 7, 7, 7, and $7 - 1 = 6$ is written back to each.  The hardware
    performs identical work for every thread of a warp; only the
    register values differ.
  \item \emph{\texttt{W3} at \texttt{0x800C}: \texttt{mov r2, \%ctaid.x}.}
    \texttt{W3} belongs to block~1, so the special register reads 1 in
    every lane and \texttt{R2} of \texttt{T4}--\texttt{T7} of block~1
    becomes 1, 1, 1, 1.  Had \texttt{W3}'s PC been \texttt{0x8008}
    instead, \texttt{mov r1, \%tid.x} would have written the thread
    indices 4, 5, 6, 7 into \texttt{R1}: the same instruction yields a
    different value per lane because the special register differs per
    thread.
\end{enumerate}
\Cref{tab:l04-walkthrough-after} shows block~0's registers after the
first two steps.  Note that \texttt{W2} has the same PC as \texttt{W0}
had: the two warps of different blocks execute the same static
instruction on their own registers, and nothing in the pipeline ties
them together.\tip{To trace a warp instruction by hand: fetch at the
warp's PC; broadcast the immediate; read the named register for every
thread of the warp; compute per lane; write back per lane; advance
that warp's PC.}

\begin{table}[htbp]
  \centering
  \caption{Register file of block~0 after steps 1 and 2 (changed values
    in bold).}
  \label{tab:l04-walkthrough-after}
  \begin{tabular}{lcccccccc}
    \toprule
    block 0 & \texttt{T0} & \texttt{T1} & \texttt{T2} & \texttt{T3} & \texttt{T4} & \texttt{T5} & \texttt{T6} & \texttt{T7} \\
    \midrule
    \texttt{R1} & \textbf{22} & \textbf{23} & \textbf{24} & \textbf{25} & 30 & 31 & 32 & 33 \\
    \texttt{R2} &  5 &  5 &  5 &  5 & \textbf{6} & \textbf{6} & \textbf{6} & \textbf{6} \\
    \texttt{R3} &  1 &  4 &  9 & 16 & 25 & 36 & 49 & 64 \\
    \texttt{R4} &  8 &  6 &  7 &  5 &  3 &  0 &  9 &  2 \\
    \bottomrule
  \end{tabular}
\end{table}

\subsection{Mask bits}

Not every instruction needs all the threads of its warp: a block of 6
threads with warp size 4 leaves two lanes of its second warp empty, and
Lesson~6 shows that a branch taken by some threads only leaves the
others idle until the paths reconverge.  The pipeline handles both with
one mechanism.\margin{The
masks of Lesson~6's SIMT stack are exactly these bits, pushed and
popped as branches diverge and reconverge; Lesson~5's warp vote and
shuffle functions expose the active mask to the program.}

\begin{definition}[Lane and active mask]
  \label{def:l04-active-mask}
  A \term*{lane}\index{lane|see{SIMT lane}}\index{SIMT lane} is one execution path through the pipeline:
  the operand-buffer entries, ALU and write-back slot that serve one
  thread of a warp instruction (the SIMT lane of Lesson~1 seen from the
  pipeline).  The \term{active mask}[アクティブマスク] of a warp is a bit
  vector with one bit per lane; a lane whose bit is 1 executes the
  instruction and writes its result back, a lane whose bit is 0 does
  no work and writes nothing.
\end{definition}

The walk-through above ran with every lane active.
\Cref{fig:l04-scoreboard-lanes} shows the per-lane operand buffers for
a later instruction under the masks of \cref{tab:l04-masks}, with
\texttt{W0} ready in all four lanes and selected, and \texttt{W1}
running under the mask \texttt{1110}: its first three
lanes receive operands and compute, its fourth lane is inactive and is
never marked ready.  A warp whose mask is all ones is the common case,
and the mask is part of the warp's state alongside its PC.

\begin{margintable}
  \caption{Active masks for a later instruction (the walk-through above ran
    with all masks \texttt{1111}); lane~3 of \texttt{W1} and lane~0 of
    \texttt{W2} are switched off.}
  \label{tab:l04-masks}
  \begin{tabular}{lc}
    \toprule
    warp & mask \\
    \midrule
    \texttt{W0} & \texttt{1111} \\
    \texttt{W1} & \texttt{1110} \\
    \texttt{W2} & \texttt{0111} \\
    \texttt{W3} & \texttt{1111} \\
    \bottomrule
  \end{tabular}
\end{margintable}

\begin{figure}[htbp]
  \centering
  % Scoreboard and operand buffer per lane (Lesson 4): four lanes, each
% holding for warps w0..w3 the two source operands (value v, ready r);
% w0 has both operands ready in every lane, is selected, and its operands
% move into the four ALUs.  Lane 3 of w1 is inactive (mask 1110) and is
% left unready.  Redraws the outline figures m4l2_scoreboard_ready,
% m4l2_scoreboard_ready_execution and m4l3_scoreboard9 (the slide's w3
% lane-1 ready mark under mask 0111 corrected).
% Usage: % Scoreboard and operand buffer per lane (Lesson 4): four lanes, each
% holding for warps w0..w3 the two source operands (value v, ready r);
% w0 has both operands ready in every lane, is selected, and its operands
% move into the four ALUs.  Lane 3 of w1 is inactive (mask 1110) and is
% left unready.  Redraws the outline figures m4l2_scoreboard_ready,
% m4l2_scoreboard_ready_execution and m4l3_scoreboard9 (the slide's w3
% lane-1 ready mark under mask 0111 corrected).
% Usage: % Scoreboard and operand buffer per lane (Lesson 4): four lanes, each
% holding for warps w0..w3 the two source operands (value v, ready r);
% w0 has both operands ready in every lane, is selected, and its operands
% move into the four ALUs.  Lane 3 of w1 is inactive (mask 1110) and is
% left unready.  Redraws the outline figures m4l2_scoreboard_ready,
% m4l2_scoreboard_ready_execution and m4l3_scoreboard9 (the slide's w3
% lane-1 ready mark under mask 0111 corrected).
% Usage: \input{figures/l04-scoreboard-lanes}
% Width: body (120 mm); the drawing is about 118 mm wide.
\begin{tikzpicture}[x=1mm, y=1mm, font=\scriptsize\sffamily,
    >={Stealth[length=1.6mm]}]
  % \lane{x}{n}{src1 value of w0}{src2 value of w0}{w1 ready marks}
  \newcommand{\lane}[5]{%
    \node[anchor=south west, font=\scriptsize\sffamily\bfseries] at (#1,30) {\iflangja{レーン #2}{lane #2}};
    % header
    \node[anchor=south] at (#1+11,25.6) {Src1};
    \node[anchor=south] at (#1+21,25.6) {Src2};
    \foreach \c/\h in {8/v, 14/r, 20/v, 26/r} {
      \node[anchor=south, text=ln-muted] at (#1+\c,23) {\h};
    }
    \foreach \w in {0,...,3} {
      \node[anchor=east, inner sep=1pt] at (#1+4,20.5-\w*4) {w\w};
      \draw[fill=ln-box] (#1+4,22.5-\w*4) rectangle ++(8,-4);
      \draw[fill=ln-box] (#1+12,22.5-\w*4) rectangle ++(4,-4);
      \draw[fill=ln-box] (#1+16,22.5-\w*4) rectangle ++(8,-4);
      \draw[fill=ln-box] (#1+24,22.5-\w*4) rectangle ++(4,-4);
    }
    % w0: both operands ready
    \node at (#1+8,20.5) {#3};
    \node[text=ln-tip] at (#1+14,20.5) {$\bullet$};
    \node at (#1+20,20.5) {#4};
    \node[text=ln-tip] at (#1+26,20.5) {$\bullet$};
    \draw[ln-tip, very thick] (#1+4,22.5) rectangle ++(24,-4);
    % w1: Src2 ready, Src1 pending (#5 = ready mark or empty)
    \node[text=ln-tip] at (#1+26,16.5) {#5};
    % ALU
    \draw[fill=ln-accent!15] (#1+6,2) -- (#1+22,2) -- (#1+18,-5) -- (#1+10,-5) -- cycle;
    \node at (#1+14,-1.5) {ALU};
    \draw[->] (#1+8,6.5) -- (#1+10,2);
    \draw[->] (#1+20,6.5) -- (#1+18,2);
  }
  \lane{0}{0}{20}{2}{$\bullet$}
  \lane{30}{1}{21}{2}{$\bullet$}
  \lane{60}{2}{22}{2}{$\bullet$}
  \lane{90}{3}{23}{2}{}
  \node[anchor=north west, align=left, text=ln-muted] at (0,-7)
    {\iflangja{w0 は全レーンで両オペランドが揃ったので選択される；w1 のレーン 3 はマスク 1110 で非アクティブ}{w0 has both operands in every lane and is selected; lane 3 of w1 is inactive under mask 1110}};
\end{tikzpicture}

% Width: body (120 mm); the drawing is about 118 mm wide.
\begin{tikzpicture}[x=1mm, y=1mm, font=\scriptsize\sffamily,
    >={Stealth[length=1.6mm]}]
  % \lane{x}{n}{src1 value of w0}{src2 value of w0}{w1 ready marks}
  \newcommand{\lane}[5]{%
    \node[anchor=south west, font=\scriptsize\sffamily\bfseries] at (#1,30) {\iflangja{レーン #2}{lane #2}};
    % header
    \node[anchor=south] at (#1+11,25.6) {Src1};
    \node[anchor=south] at (#1+21,25.6) {Src2};
    \foreach \c/\h in {8/v, 14/r, 20/v, 26/r} {
      \node[anchor=south, text=ln-muted] at (#1+\c,23) {\h};
    }
    \foreach \w in {0,...,3} {
      \node[anchor=east, inner sep=1pt] at (#1+4,20.5-\w*4) {w\w};
      \draw[fill=ln-box] (#1+4,22.5-\w*4) rectangle ++(8,-4);
      \draw[fill=ln-box] (#1+12,22.5-\w*4) rectangle ++(4,-4);
      \draw[fill=ln-box] (#1+16,22.5-\w*4) rectangle ++(8,-4);
      \draw[fill=ln-box] (#1+24,22.5-\w*4) rectangle ++(4,-4);
    }
    % w0: both operands ready
    \node at (#1+8,20.5) {#3};
    \node[text=ln-tip] at (#1+14,20.5) {$\bullet$};
    \node at (#1+20,20.5) {#4};
    \node[text=ln-tip] at (#1+26,20.5) {$\bullet$};
    \draw[ln-tip, very thick] (#1+4,22.5) rectangle ++(24,-4);
    % w1: Src2 ready, Src1 pending (#5 = ready mark or empty)
    \node[text=ln-tip] at (#1+26,16.5) {#5};
    % ALU
    \draw[fill=ln-accent!15] (#1+6,2) -- (#1+22,2) -- (#1+18,-5) -- (#1+10,-5) -- cycle;
    \node at (#1+14,-1.5) {ALU};
    \draw[->] (#1+8,6.5) -- (#1+10,2);
    \draw[->] (#1+20,6.5) -- (#1+18,2);
  }
  \lane{0}{0}{20}{2}{$\bullet$}
  \lane{30}{1}{21}{2}{$\bullet$}
  \lane{60}{2}{22}{2}{$\bullet$}
  \lane{90}{3}{23}{2}{}
  \node[anchor=north west, align=left, text=ln-muted] at (0,-7)
    {\iflangja{w0 は全レーンで両オペランドが揃ったので選択される；w1 のレーン 3 はマスク 1110 で非アクティブ}{w0 has both operands in every lane and is selected; lane 3 of w1 is inactive under mask 1110}};
\end{tikzpicture}

% Width: body (120 mm); the drawing is about 118 mm wide.
\begin{tikzpicture}[x=1mm, y=1mm, font=\scriptsize\sffamily,
    >={Stealth[length=1.6mm]}]
  % \lane{x}{n}{src1 value of w0}{src2 value of w0}{w1 ready marks}
  \newcommand{\lane}[5]{%
    \node[anchor=south west, font=\scriptsize\sffamily\bfseries] at (#1,30) {\iflangja{レーン #2}{lane #2}};
    % header
    \node[anchor=south] at (#1+11,25.6) {Src1};
    \node[anchor=south] at (#1+21,25.6) {Src2};
    \foreach \c/\h in {8/v, 14/r, 20/v, 26/r} {
      \node[anchor=south, text=ln-muted] at (#1+\c,23) {\h};
    }
    \foreach \w in {0,...,3} {
      \node[anchor=east, inner sep=1pt] at (#1+4,20.5-\w*4) {w\w};
      \draw[fill=ln-box] (#1+4,22.5-\w*4) rectangle ++(8,-4);
      \draw[fill=ln-box] (#1+12,22.5-\w*4) rectangle ++(4,-4);
      \draw[fill=ln-box] (#1+16,22.5-\w*4) rectangle ++(8,-4);
      \draw[fill=ln-box] (#1+24,22.5-\w*4) rectangle ++(4,-4);
    }
    % w0: both operands ready
    \node at (#1+8,20.5) {#3};
    \node[text=ln-tip] at (#1+14,20.5) {$\bullet$};
    \node at (#1+20,20.5) {#4};
    \node[text=ln-tip] at (#1+26,20.5) {$\bullet$};
    \draw[ln-tip, very thick] (#1+4,22.5) rectangle ++(24,-4);
    % w1: Src2 ready, Src1 pending (#5 = ready mark or empty)
    \node[text=ln-tip] at (#1+26,16.5) {#5};
    % ALU
    \draw[fill=ln-accent!15] (#1+6,2) -- (#1+22,2) -- (#1+18,-5) -- (#1+10,-5) -- cycle;
    \node at (#1+14,-1.5) {ALU};
    \draw[->] (#1+8,6.5) -- (#1+10,2);
    \draw[->] (#1+20,6.5) -- (#1+18,2);
  }
  \lane{0}{0}{20}{2}{$\bullet$}
  \lane{30}{1}{21}{2}{$\bullet$}
  \lane{60}{2}{22}{2}{$\bullet$}
  \lane{90}{3}{23}{2}{}
  \node[anchor=north west, align=left, text=ln-muted] at (0,-7)
    {\iflangja{w0 は全レーンで両オペランドが揃ったので選択される；w1 のレーン 3 はマスク 1110 で非アクティブ}{w0 has both operands in every lane and is selected; lane 3 of w1 is inactive under mask 1110}};
\end{tikzpicture}

  \caption{Per-lane operand buffers and the scoreboard: \texttt{w0}
    has both operands ready in every lane and is issued to the four
    ALUs; \texttt{w1} runs under the mask \texttt{1110}, so its fourth
    lane stays inactive.}
  \label{fig:l04-scoreboard-lanes}
\end{figure}


\begin{keypoint}
  A warp instruction is fetched once at the warp's PC, decoded once,
  and then read, executed and written back per lane; immediates and
  special registers like \texttt{\%ctaid.x} are broadcast, while
  \texttt{\%tid.x} and ordinary registers differ per lane.  The active
  mask switches individual lanes off without changing the instruction
  stream.
\end{keypoint}

\section{Global memory coalescing}
\label{sec:l04-coalescing}

The execution side of the SM processes a warp instruction as one unit;
the memory side cannot, because each thread of a load supplies its own
address.\source{M4, coalescing; K12 §1.4.3; A18 ch.~4; BPG §10.2.1
Coalesced Access to Global Memory; CPG §5.3.2 Device Memory Accesses;
PMPP ch.~6.}  This section counts the requests a warp can generate,
reviews the memory technologies that must serve them, and shows how the
hardware merges requests whose addresses are close.

\subsection{The demand for bandwidth}

One load instruction of a 32-thread warp can generate up to 32 distinct
memory requests, one per thread.  The \term{memory bandwidth}[メモリ帯域幅]
that would be needed if every SM issued such a load every cycle is
easily computed.

\begin{example}
  \label{ex:l04-bandwidth-demand}
  A GPU has 64 SMs at \SI{1.5}{\giga\hertz}.  If every SM issues one
  warp load per cycle and every thread's request fetches
  \SI{32}{\byte}, the SMs demand
  $64 \times 32 \times \SI{32}{\byte} = \SI{65536}{\byte}$ per cycle,
  which at \SI{1.5}{\giga\hertz} is about
  \SI{98}{\tera\byte\per\second}.  (With 32 SMs, \SI{64}{\byte} per
  request and \SI{1}{\giga\hertz} the same \SI{64}{\kibi\byte} per cycle
  give about \SI{66}{\tera\byte\per\second}, usually rounded to
  \SI{64}{\tera\byte\per\second}.)  The H100 SXM5 memory system delivers
  about \SI{3}{\tera\byte\per\second} (preliminary figure), some thirty
  times less.\margin{HBM bandwidth by generation: P100
  \SI{720}{\giga\byte\per\second}, V100 \SI{900}{\giga\byte\per\second},
  A100 (\SI{40}{\gibi\byte}) \SI{1555}{\giga\byte\per\second}
  \cite{nvidia2020ampere}; H100 SXM5 \SI{3000}{\giga\byte\per\second}
  (preliminary, HWP V1.02) \cite{nvidia2022hopper}.}  The gap is closed
  by caches, by not every warp loading every cycle, and above all by
  coalescing: the 32 requests of a warp usually fall into a few
  adjacent blocks of memory.
\end{example}

GPU caches are small relative to the thread count, so their bandwidth
saturates easily, and the off-chip memory is the bottleneck for most
kernels.  Which memory technology sits there decides how wide the
bottleneck is.

\subsection{SRAM, DRAM and HBM}

\Cref{fig:l04-sram-dram-hbm} contrasts the three technologies.

\begin{description}
  \item[SRAM.]  Static random-access memory, \term{SRAM}[SRAM], stores a
    bit in a six-transistor cell: two cross-coupled inverters that hold
    the bit, connected to the bit lines through two access transistors
    gated by the word line.  It is fast and needs no refresh, but
    six transistors per bit make it expensive in area, so it is used
    for the structures on the GPU die: registers, shared memory and
    caches.
  \item[DRAM.]  Dynamic random-access memory, \term{DRAM}[DRAM], stores a
    bit as charge on a capacitor behind a single transistor (1T1C).
    The cell is far smaller, so DRAM gives the capacity of the GPU's
    main memory, but it is slower, needs periodic refresh, and its chips
    sit in separate packages: every bit that enters or leaves must pass
    through the package pins, whose number and speed bound the
    bandwidth.  The SMs reach the L2 and DRAM through several
    \emph{memory channels}, each with its own DRAM, and coalesced
    requests spread over them (K12 §1.4.3 \cite{kim2012gpgpu}).
  \item[HBM.]  High bandwidth memory (Lesson~1) keeps the DRAM cell but
    changes the packaging.  Several DRAM dies are stacked vertically and
    connected through the stack, which is
    \term{3D-stacked memory}[3次元積層メモリ], and the stack is placed
    next to the GPU die on a \term{silicon interposer}[シリコンインターポーザ],
    a silicon substrate carrying thousands of fine wires inside the same
    package.  The link to the GPU no longer crosses package pins or a
    circuit board, so it can be thousands of bits wide, and the bandwidth
    per package is several times that of a GDDR memory system.\margin{Earlier GPUs used GDDR DRAM over a wide bus, with an FR-FCFS
    (first-ready, first-come first-served) scheduler at each channel: K12
    §1.4.3 \cite{kim2012gpgpu}.  Interface widths:
    P100 and V100 use a 4096-bit HBM2 interface, A100 a 5120-bit HBM2
    interface \cite{nvidia2020ampere}, H100 SXM5 a 5120-bit HBM3
    interface \cite{nvidia2022hopper}.  BPG §9.2.1 derives V100's
    \SI{898}{\giga\byte\per\second} as $0.877\,\text{GHz} \times 2
    \times 4096\,\text{bit} / 8$ \cite{nvidia2026cudabestpractices}.}
\end{description}

\begin{figure}[htbp]
  \centering
  % SRAM, DRAM and HBM (Lesson 4).  (a) a six-transistor SRAM cell: two
% cross-coupled inverters and two access transistors on a word line and a
% pair of bit lines.  (b) a one-transistor, one-capacitor DRAM cell, and
% a DRAM package whose chips talk through pins.  (c) a GPU package with a
% silicon interposer carrying the GPU die and a 3D stack of DRAM dies.
% Schematic symbols only.  Redraws the outline figure m4l4_sram.
% Usage: % SRAM, DRAM and HBM (Lesson 4).  (a) a six-transistor SRAM cell: two
% cross-coupled inverters and two access transistors on a word line and a
% pair of bit lines.  (b) a one-transistor, one-capacitor DRAM cell, and
% a DRAM package whose chips talk through pins.  (c) a GPU package with a
% silicon interposer carrying the GPU die and a 3D stack of DRAM dies.
% Schematic symbols only.  Redraws the outline figure m4l4_sram.
% Usage: % SRAM, DRAM and HBM (Lesson 4).  (a) a six-transistor SRAM cell: two
% cross-coupled inverters and two access transistors on a word line and a
% pair of bit lines.  (b) a one-transistor, one-capacitor DRAM cell, and
% a DRAM package whose chips talk through pins.  (c) a GPU package with a
% silicon interposer carrying the GPU die and a 3D stack of DRAM dies.
% Schematic symbols only.  Redraws the outline figure m4l4_sram.
% Usage: \input{figures/l04-sram-dram-hbm}
% Width: body (120 mm); the drawing is about 118 mm wide.
\begin{tikzpicture}[x=1mm, y=1mm, font=\scriptsize\sffamily,
    >={Stealth[length=1.6mm]}]
  % (a) SRAM cell
  \node[anchor=south west, font=\scriptsize\sffamily\bfseries] at (0,33.5)
    {\iflangja{(a) SRAM セル（6 トランジスタ）}{(a) SRAM cell (6 transistors)}};
  \draw[ln-muted] (2,25) -- (34,25) node[right, text=ln-muted] {\iflangja{ワード線}{word line}};
  \draw[ln-muted] (4,4) -- (4,28.5) node[above, text=ln-muted] {\iflangja{ビット線}{bit line}};
  \draw[ln-muted] (32,4) -- (32,28.5) node[above, text=ln-muted] {\iflangja{$\overline{\text{ビット線}}$}{$\overline{\text{bit line}}$}};
  % two inverters, cross coupled: the upper one drives Q-bar from Q, the
  % lower one drives Q from Q-bar; storage nodes Q (x=10) and Q-bar (x=26)
  \draw[fill=ln-box] (14,22) -- (14,16) -- (21,19) -- cycle;
  \draw[fill=white] (21.9,19) circle (0.9);
  \draw (10,19) -- (14,19); \draw (22.8,19) -- (26,19);
  \draw[fill=ln-box] (22,14) -- (22,8) -- (15,11) -- cycle;
  \draw[fill=white] (14.1,11) circle (0.9);
  \draw (26,11) -- (22,11); \draw (13.2,11) -- (10,11);
  \draw (10,11) -- (10,19); \draw (26,11) -- (26,19);
  \fill (10,15) circle (0.45); \fill (26,15) circle (0.45);
  \node[anchor=north] at (10,10.6) {$Q$};
  \node[anchor=north] at (26,10.6) {$\overline{Q}$};
  % access transistors (schematic boxes): bit line to Q, bit line-bar to
  % Q-bar, both gated by the word line
  \draw[fill=ln-accent!15] (5,13) rectangle ++(4,4);
  \draw (4,15) -- (5,15); \draw (9,15) -- (10,15); \draw (7,17) -- (7,25);
  \draw[fill=ln-accent!15] (27,13) rectangle ++(4,4);
  \draw (31,15) -- (32,15); \draw (26,15) -- (27,15); \draw (29,17) -- (29,25);
  \node[anchor=north west, text=ln-muted, align=left] at (0,3)
    {\iflangja{キャッシュ，レジスタ，共有メモリ}{caches, registers, shared memory}};
  % (b) DRAM cell and package
  \node[anchor=south west, font=\scriptsize\sffamily\bfseries] at (44,33.5)
    {\iflangja{(b) DRAM（1T1C）とパッケージ}{(b) DRAM cell (1T1C) and package}};
  \draw[ln-muted] (46,25) -- (66,25) node[right, text=ln-muted] {\iflangja{ワード線}{word line}};
  \draw[ln-muted] (48,13) -- (48,28.5) node[above, text=ln-muted] {\iflangja{ビット線}{bit line}};
  \draw[fill=ln-accent!15] (51,18) rectangle ++(4,4);
  \draw (48,20) -- (51,20); \draw (55,20) -- (58,20); \draw (53,22) -- (53,25);
  \draw (58,20) -- (58,18.5); \draw (56,18.5) -- (60,18.5); \draw (56,17) -- (60,17);
  \draw (58,17) -- (58,15.5); \draw (56.5,15.5) -- (59.5,15.5);
  \node[anchor=west, text=ln-muted] at (61,17.5) {\iflangja{キャパシタ}{capacitor}};
  % package with chips and pins
  \draw[fill=ln-box] (46,2) rectangle (76,10);
  \foreach \k in {0,...,3} { \draw[fill=ln-tip!25] (47.5+\k*7.2,4) rectangle ++(6,4); }
  \foreach \k in {0,...,11} { \draw (47+\k*2.5,2) -- ++(0,-2); }
  \node[anchor=north west, text=ln-muted] at (44,-0.5) {\iflangja{全通信がピンを通る}{all traffic through the pins}};
  % (c) HBM package
  \node[anchor=south west, font=\scriptsize\sffamily\bfseries] at (84,33.5)
    {\iflangja{(c) HBM の GPU パッケージ}{(c) GPU package with HBM}};
  \draw[fill=ln-box] (84,2) rectangle (118,26);
  \draw[fill=ln-accent!15] (86,6) rectangle (116,10)
    node[midway] {\iflangja{シリコンインターポーザ}{silicon interposer}};
  \draw[fill=ln-accent!25] (87,10) rectangle (101,22) node[midway] {GPU};
  \foreach \k in {0,...,4} {
    \draw[fill=ln-tip!25] (104,10+\k*1.8) rectangle ++(11,1.5);
  }
  \node[anchor=south, align=center] at (109.5,19.0) {\iflangja{3D 積層\\DRAM}{3D-stacked\\DRAM}};
  \draw[<->, thick, ln-caution] (101,12) -- (104,12);
  \node[anchor=north west, text=ln-muted, align=left] at (84,1)
    {\iflangja{数千本の短い配線}{thousands of short wires}};
\end{tikzpicture}

% Width: body (120 mm); the drawing is about 118 mm wide.
\begin{tikzpicture}[x=1mm, y=1mm, font=\scriptsize\sffamily,
    >={Stealth[length=1.6mm]}]
  % (a) SRAM cell
  \node[anchor=south west, font=\scriptsize\sffamily\bfseries] at (0,33.5)
    {\iflangja{(a) SRAM セル（6 トランジスタ）}{(a) SRAM cell (6 transistors)}};
  \draw[ln-muted] (2,25) -- (34,25) node[right, text=ln-muted] {\iflangja{ワード線}{word line}};
  \draw[ln-muted] (4,4) -- (4,28.5) node[above, text=ln-muted] {\iflangja{ビット線}{bit line}};
  \draw[ln-muted] (32,4) -- (32,28.5) node[above, text=ln-muted] {\iflangja{$\overline{\text{ビット線}}$}{$\overline{\text{bit line}}$}};
  % two inverters, cross coupled: the upper one drives Q-bar from Q, the
  % lower one drives Q from Q-bar; storage nodes Q (x=10) and Q-bar (x=26)
  \draw[fill=ln-box] (14,22) -- (14,16) -- (21,19) -- cycle;
  \draw[fill=white] (21.9,19) circle (0.9);
  \draw (10,19) -- (14,19); \draw (22.8,19) -- (26,19);
  \draw[fill=ln-box] (22,14) -- (22,8) -- (15,11) -- cycle;
  \draw[fill=white] (14.1,11) circle (0.9);
  \draw (26,11) -- (22,11); \draw (13.2,11) -- (10,11);
  \draw (10,11) -- (10,19); \draw (26,11) -- (26,19);
  \fill (10,15) circle (0.45); \fill (26,15) circle (0.45);
  \node[anchor=north] at (10,10.6) {$Q$};
  \node[anchor=north] at (26,10.6) {$\overline{Q}$};
  % access transistors (schematic boxes): bit line to Q, bit line-bar to
  % Q-bar, both gated by the word line
  \draw[fill=ln-accent!15] (5,13) rectangle ++(4,4);
  \draw (4,15) -- (5,15); \draw (9,15) -- (10,15); \draw (7,17) -- (7,25);
  \draw[fill=ln-accent!15] (27,13) rectangle ++(4,4);
  \draw (31,15) -- (32,15); \draw (26,15) -- (27,15); \draw (29,17) -- (29,25);
  \node[anchor=north west, text=ln-muted, align=left] at (0,3)
    {\iflangja{キャッシュ，レジスタ，共有メモリ}{caches, registers, shared memory}};
  % (b) DRAM cell and package
  \node[anchor=south west, font=\scriptsize\sffamily\bfseries] at (44,33.5)
    {\iflangja{(b) DRAM（1T1C）とパッケージ}{(b) DRAM cell (1T1C) and package}};
  \draw[ln-muted] (46,25) -- (66,25) node[right, text=ln-muted] {\iflangja{ワード線}{word line}};
  \draw[ln-muted] (48,13) -- (48,28.5) node[above, text=ln-muted] {\iflangja{ビット線}{bit line}};
  \draw[fill=ln-accent!15] (51,18) rectangle ++(4,4);
  \draw (48,20) -- (51,20); \draw (55,20) -- (58,20); \draw (53,22) -- (53,25);
  \draw (58,20) -- (58,18.5); \draw (56,18.5) -- (60,18.5); \draw (56,17) -- (60,17);
  \draw (58,17) -- (58,15.5); \draw (56.5,15.5) -- (59.5,15.5);
  \node[anchor=west, text=ln-muted] at (61,17.5) {\iflangja{キャパシタ}{capacitor}};
  % package with chips and pins
  \draw[fill=ln-box] (46,2) rectangle (76,10);
  \foreach \k in {0,...,3} { \draw[fill=ln-tip!25] (47.5+\k*7.2,4) rectangle ++(6,4); }
  \foreach \k in {0,...,11} { \draw (47+\k*2.5,2) -- ++(0,-2); }
  \node[anchor=north west, text=ln-muted] at (44,-0.5) {\iflangja{全通信がピンを通る}{all traffic through the pins}};
  % (c) HBM package
  \node[anchor=south west, font=\scriptsize\sffamily\bfseries] at (84,33.5)
    {\iflangja{(c) HBM の GPU パッケージ}{(c) GPU package with HBM}};
  \draw[fill=ln-box] (84,2) rectangle (118,26);
  \draw[fill=ln-accent!15] (86,6) rectangle (116,10)
    node[midway] {\iflangja{シリコンインターポーザ}{silicon interposer}};
  \draw[fill=ln-accent!25] (87,10) rectangle (101,22) node[midway] {GPU};
  \foreach \k in {0,...,4} {
    \draw[fill=ln-tip!25] (104,10+\k*1.8) rectangle ++(11,1.5);
  }
  \node[anchor=south, align=center] at (109.5,19.0) {\iflangja{3D 積層\\DRAM}{3D-stacked\\DRAM}};
  \draw[<->, thick, ln-caution] (101,12) -- (104,12);
  \node[anchor=north west, text=ln-muted, align=left] at (84,1)
    {\iflangja{数千本の短い配線}{thousands of short wires}};
\end{tikzpicture}

% Width: body (120 mm); the drawing is about 118 mm wide.
\begin{tikzpicture}[x=1mm, y=1mm, font=\scriptsize\sffamily,
    >={Stealth[length=1.6mm]}]
  % (a) SRAM cell
  \node[anchor=south west, font=\scriptsize\sffamily\bfseries] at (0,33.5)
    {\iflangja{(a) SRAM セル（6 トランジスタ）}{(a) SRAM cell (6 transistors)}};
  \draw[ln-muted] (2,25) -- (34,25) node[right, text=ln-muted] {\iflangja{ワード線}{word line}};
  \draw[ln-muted] (4,4) -- (4,28.5) node[above, text=ln-muted] {\iflangja{ビット線}{bit line}};
  \draw[ln-muted] (32,4) -- (32,28.5) node[above, text=ln-muted] {\iflangja{$\overline{\text{ビット線}}$}{$\overline{\text{bit line}}$}};
  % two inverters, cross coupled: the upper one drives Q-bar from Q, the
  % lower one drives Q from Q-bar; storage nodes Q (x=10) and Q-bar (x=26)
  \draw[fill=ln-box] (14,22) -- (14,16) -- (21,19) -- cycle;
  \draw[fill=white] (21.9,19) circle (0.9);
  \draw (10,19) -- (14,19); \draw (22.8,19) -- (26,19);
  \draw[fill=ln-box] (22,14) -- (22,8) -- (15,11) -- cycle;
  \draw[fill=white] (14.1,11) circle (0.9);
  \draw (26,11) -- (22,11); \draw (13.2,11) -- (10,11);
  \draw (10,11) -- (10,19); \draw (26,11) -- (26,19);
  \fill (10,15) circle (0.45); \fill (26,15) circle (0.45);
  \node[anchor=north] at (10,10.6) {$Q$};
  \node[anchor=north] at (26,10.6) {$\overline{Q}$};
  % access transistors (schematic boxes): bit line to Q, bit line-bar to
  % Q-bar, both gated by the word line
  \draw[fill=ln-accent!15] (5,13) rectangle ++(4,4);
  \draw (4,15) -- (5,15); \draw (9,15) -- (10,15); \draw (7,17) -- (7,25);
  \draw[fill=ln-accent!15] (27,13) rectangle ++(4,4);
  \draw (31,15) -- (32,15); \draw (26,15) -- (27,15); \draw (29,17) -- (29,25);
  \node[anchor=north west, text=ln-muted, align=left] at (0,3)
    {\iflangja{キャッシュ，レジスタ，共有メモリ}{caches, registers, shared memory}};
  % (b) DRAM cell and package
  \node[anchor=south west, font=\scriptsize\sffamily\bfseries] at (44,33.5)
    {\iflangja{(b) DRAM（1T1C）とパッケージ}{(b) DRAM cell (1T1C) and package}};
  \draw[ln-muted] (46,25) -- (66,25) node[right, text=ln-muted] {\iflangja{ワード線}{word line}};
  \draw[ln-muted] (48,13) -- (48,28.5) node[above, text=ln-muted] {\iflangja{ビット線}{bit line}};
  \draw[fill=ln-accent!15] (51,18) rectangle ++(4,4);
  \draw (48,20) -- (51,20); \draw (55,20) -- (58,20); \draw (53,22) -- (53,25);
  \draw (58,20) -- (58,18.5); \draw (56,18.5) -- (60,18.5); \draw (56,17) -- (60,17);
  \draw (58,17) -- (58,15.5); \draw (56.5,15.5) -- (59.5,15.5);
  \node[anchor=west, text=ln-muted] at (61,17.5) {\iflangja{キャパシタ}{capacitor}};
  % package with chips and pins
  \draw[fill=ln-box] (46,2) rectangle (76,10);
  \foreach \k in {0,...,3} { \draw[fill=ln-tip!25] (47.5+\k*7.2,4) rectangle ++(6,4); }
  \foreach \k in {0,...,11} { \draw (47+\k*2.5,2) -- ++(0,-2); }
  \node[anchor=north west, text=ln-muted] at (44,-0.5) {\iflangja{全通信がピンを通る}{all traffic through the pins}};
  % (c) HBM package
  \node[anchor=south west, font=\scriptsize\sffamily\bfseries] at (84,33.5)
    {\iflangja{(c) HBM の GPU パッケージ}{(c) GPU package with HBM}};
  \draw[fill=ln-box] (84,2) rectangle (118,26);
  \draw[fill=ln-accent!15] (86,6) rectangle (116,10)
    node[midway] {\iflangja{シリコンインターポーザ}{silicon interposer}};
  \draw[fill=ln-accent!25] (87,10) rectangle (101,22) node[midway] {GPU};
  \foreach \k in {0,...,4} {
    \draw[fill=ln-tip!25] (104,10+\k*1.8) rectangle ++(11,1.5);
  }
  \node[anchor=south, align=center] at (109.5,19.0) {\iflangja{3D 積層\\DRAM}{3D-stacked\\DRAM}};
  \draw[<->, thick, ln-caution] (101,12) -- (104,12);
  \node[anchor=north west, text=ln-muted, align=left] at (84,1)
    {\iflangja{数千本の短い配線}{thousands of short wires}};
\end{tikzpicture}

  \caption{(a) A six-transistor SRAM cell.  (b) A one-transistor,
    one-capacitor DRAM cell and a DRAM package whose traffic passes
    through its pins.  (c) A GPU package in which a silicon interposer
    links the GPU die to a 3D stack of DRAM dies.}
  \label{fig:l04-sram-dram-hbm}
\end{figure}

\subsection{Coalescing}

Even HBM cannot serve 32 separate requests per warp instruction, so the
memory unit of the SM looks at the 32 addresses together.

\begin{definition}[Memory coalescing]
  \label{def:l04-coalescing}
  The combination, by the hardware, of the memory requests of the
  threads of one warp instruction into as few memory transactions as
  their addresses allow is \term{memory coalescing}[メモリコアレッシング].
  A \term{memory transaction}[メモリトランザクション] is one transfer
  between the memory system and the SM of a fixed, naturally aligned
  block of bytes.  An access whose 32 requests fit into a small number
  of transactions is a \term{coalesced access}[コアレスアクセス]; one that
  needs a separate transaction for most threads is an
  \term{uncoalesced access}[非コアレスアクセス].
\end{definition}

\Cref{fig:l04-coalescing} shows the two cases for eight threads
executing two PTX loads,
\begin{lstlisting}[language=PTX, numbers=none]
ld.global.u32 %r3, [%r1];   // R1: 512, 516, ..., 540
ld.global.u32 %r3, [%r2];   // R2: 0, 128, 256, ..., 896
\end{lstlisting}
where a DRAM request is taken to be \SI{64}{\byte} long.\caution{Lecture vs guide: the \SI{64}{\byte} request of the
example is the granularity the lecture assumes (\enquote{64 to
128~B}).  Since compute capability 6.0 the hardware counts
\SI{32}{\byte} sectors: BPG §10.2.1 states that a warp's accesses
coalesce into \enquote{a number of transactions equal to the number of
32-byte transactions necessary to service all of the threads of the
warp} \cite{nvidia2026cudabestpractices}.  Lesson~8 returns to sectored
caches.}  In the first
load the eight addresses are consecutive 4-byte words from 512 to 540,
\SI{32}{\byte} in all, so they combine into one request for
\texttt{Mem[512--575]}: the access is coalesced.  In the second the
addresses are \SI{128}{\byte} apart, so no two of them share a request
and eight requests are issued: the access is uncoalesced and consumes
eight times the bandwidth for the same \SI{32}{\byte} of useful
data.

\begin{figure}[htbp]
  \centering
  % Coalesced and uncoalesced loads (Lesson 4).  Eight threads execute
% ld.global.u32 %r3, [%r1] and ld.global.u32 %r3, [%r2].  R1 holds eight
% consecutive words, so the eight accesses combine into one 64-byte
% request; R2 holds addresses 128 B apart, so eight requests are needed.
% Redraws the outline figures m4l4_coalesce and m4l4_coalesce2 (the
% slide's address typos corrected; a uniform stride is used).
% Usage: % Coalesced and uncoalesced loads (Lesson 4).  Eight threads execute
% ld.global.u32 %r3, [%r1] and ld.global.u32 %r3, [%r2].  R1 holds eight
% consecutive words, so the eight accesses combine into one 64-byte
% request; R2 holds addresses 128 B apart, so eight requests are needed.
% Redraws the outline figures m4l4_coalesce and m4l4_coalesce2 (the
% slide's address typos corrected; a uniform stride is used).
% Usage: % Coalesced and uncoalesced loads (Lesson 4).  Eight threads execute
% ld.global.u32 %r3, [%r1] and ld.global.u32 %r3, [%r2].  R1 holds eight
% consecutive words, so the eight accesses combine into one 64-byte
% request; R2 holds addresses 128 B apart, so eight requests are needed.
% Redraws the outline figures m4l4_coalesce and m4l4_coalesce2 (the
% slide's address typos corrected; a uniform stride is used).
% Usage: \input{figures/l04-coalescing}
% Width: body (120 mm); the drawing is about 118 mm wide.
\begin{tikzpicture}[x=1mm, y=1mm, font=\scriptsize\sffamily,
    >={Stealth[length=1.6mm]}]
  % register table
  \node[anchor=south west, font=\scriptsize\sffamily\bfseries] at (0,41)
    {\iflangja{レジスタ}{registers}};
  \node[anchor=south] at (16,35.5) {R1};
  \node[anchor=south] at (28,35.5) {R2};
  \foreach \t/\a/\b [count=\i from 0] in {0/512/0, 1/516/128, 2/520/256, 3/524/384, 4/528/512, 5/532/640, 6/536/768, 7/540/896} {
    \node[anchor=east, inner sep=1pt] at (9.5,33-\i*4.2) {T\t};
    \draw[fill=ln-tip!20] (10,35-\i*4.2) rectangle ++(12,-4) node[midway, font=\scriptsize\ttfamily] {\a};
    \draw[fill=ln-caution!15] (22,35-\i*4.2) rectangle ++(12,-4) node[midway, font=\scriptsize\ttfamily] {\b};
  }
  % one request for R1
  \node[anchor=south west, font=\scriptsize\sffamily\bfseries] at (46,41)
    {\iflangja{ld [\%r1]：1 リクエスト}{ld [\%r1]: one request}};
  \draw[fill=ln-tip!20, ln-tip, thick] (46,27) rectangle ++(26,-8)
    node[midway, align=center, text=black] {Mem[512--575]};
  \draw[->, ln-tip] (22,18) -- (46,23);
  \node[anchor=north west, align=left, text=ln-muted] at (46,17)
    {\iflangja{連続する 8 語（32 B）が\\1 回の 64 B リクエストに}{8 consecutive words (32 B)\\in one 64-B request}};
  % eight requests for R2
  \node[anchor=south west, font=\scriptsize\sffamily\bfseries] at (84,41)
    {\iflangja{ld [\%r2]：8 リクエスト}{ld [\%r2]: eight requests}};
  \foreach \b [count=\i from 0] in {0, 128, 256, 384, 512, 640, 768, 896} {
    \pgfmathtruncatemacro{\e}{\b+63}
    \draw[fill=ln-caution!15, ln-caution] (84,35-\i*4.2) rectangle ++(30,-3.6)
      node[midway, text=black] {Mem[\b--\e]};
  }
  \draw[->, ln-caution] (34,18) -- (84,18);
\end{tikzpicture}

% Width: body (120 mm); the drawing is about 118 mm wide.
\begin{tikzpicture}[x=1mm, y=1mm, font=\scriptsize\sffamily,
    >={Stealth[length=1.6mm]}]
  % register table
  \node[anchor=south west, font=\scriptsize\sffamily\bfseries] at (0,41)
    {\iflangja{レジスタ}{registers}};
  \node[anchor=south] at (16,35.5) {R1};
  \node[anchor=south] at (28,35.5) {R2};
  \foreach \t/\a/\b [count=\i from 0] in {0/512/0, 1/516/128, 2/520/256, 3/524/384, 4/528/512, 5/532/640, 6/536/768, 7/540/896} {
    \node[anchor=east, inner sep=1pt] at (9.5,33-\i*4.2) {T\t};
    \draw[fill=ln-tip!20] (10,35-\i*4.2) rectangle ++(12,-4) node[midway, font=\scriptsize\ttfamily] {\a};
    \draw[fill=ln-caution!15] (22,35-\i*4.2) rectangle ++(12,-4) node[midway, font=\scriptsize\ttfamily] {\b};
  }
  % one request for R1
  \node[anchor=south west, font=\scriptsize\sffamily\bfseries] at (46,41)
    {\iflangja{ld [\%r1]：1 リクエスト}{ld [\%r1]: one request}};
  \draw[fill=ln-tip!20, ln-tip, thick] (46,27) rectangle ++(26,-8)
    node[midway, align=center, text=black] {Mem[512--575]};
  \draw[->, ln-tip] (22,18) -- (46,23);
  \node[anchor=north west, align=left, text=ln-muted] at (46,17)
    {\iflangja{連続する 8 語（32 B）が\\1 回の 64 B リクエストに}{8 consecutive words (32 B)\\in one 64-B request}};
  % eight requests for R2
  \node[anchor=south west, font=\scriptsize\sffamily\bfseries] at (84,41)
    {\iflangja{ld [\%r2]：8 リクエスト}{ld [\%r2]: eight requests}};
  \foreach \b [count=\i from 0] in {0, 128, 256, 384, 512, 640, 768, 896} {
    \pgfmathtruncatemacro{\e}{\b+63}
    \draw[fill=ln-caution!15, ln-caution] (84,35-\i*4.2) rectangle ++(30,-3.6)
      node[midway, text=black] {Mem[\b--\e]};
  }
  \draw[->, ln-caution] (34,18) -- (84,18);
\end{tikzpicture}

% Width: body (120 mm); the drawing is about 118 mm wide.
\begin{tikzpicture}[x=1mm, y=1mm, font=\scriptsize\sffamily,
    >={Stealth[length=1.6mm]}]
  % register table
  \node[anchor=south west, font=\scriptsize\sffamily\bfseries] at (0,41)
    {\iflangja{レジスタ}{registers}};
  \node[anchor=south] at (16,35.5) {R1};
  \node[anchor=south] at (28,35.5) {R2};
  \foreach \t/\a/\b [count=\i from 0] in {0/512/0, 1/516/128, 2/520/256, 3/524/384, 4/528/512, 5/532/640, 6/536/768, 7/540/896} {
    \node[anchor=east, inner sep=1pt] at (9.5,33-\i*4.2) {T\t};
    \draw[fill=ln-tip!20] (10,35-\i*4.2) rectangle ++(12,-4) node[midway, font=\scriptsize\ttfamily] {\a};
    \draw[fill=ln-caution!15] (22,35-\i*4.2) rectangle ++(12,-4) node[midway, font=\scriptsize\ttfamily] {\b};
  }
  % one request for R1
  \node[anchor=south west, font=\scriptsize\sffamily\bfseries] at (46,41)
    {\iflangja{ld [\%r1]：1 リクエスト}{ld [\%r1]: one request}};
  \draw[fill=ln-tip!20, ln-tip, thick] (46,27) rectangle ++(26,-8)
    node[midway, align=center, text=black] {Mem[512--575]};
  \draw[->, ln-tip] (22,18) -- (46,23);
  \node[anchor=north west, align=left, text=ln-muted] at (46,17)
    {\iflangja{連続する 8 語（32 B）が\\1 回の 64 B リクエストに}{8 consecutive words (32 B)\\in one 64-B request}};
  % eight requests for R2
  \node[anchor=south west, font=\scriptsize\sffamily\bfseries] at (84,41)
    {\iflangja{ld [\%r2]：8 リクエスト}{ld [\%r2]: eight requests}};
  \foreach \b [count=\i from 0] in {0, 128, 256, 384, 512, 640, 768, 896} {
    \pgfmathtruncatemacro{\e}{\b+63}
    \draw[fill=ln-caution!15, ln-caution] (84,35-\i*4.2) rectangle ++(30,-3.6)
      node[midway, text=black] {Mem[\b--\e]};
  }
  \draw[->, ln-caution] (34,18) -- (84,18);
\end{tikzpicture}

  \caption{Eight threads load through \texttt{\%r1} (consecutive words,
    one request) and through \texttt{\%r2} (words \SI{128}{\byte} apart,
    eight requests).}
  \label{fig:l04-coalescing}
\end{figure}

The precise rule has three ingredients.  Global memory is accessed in
transactions of 32, 64 or \SI{128}{\byte}, and each transaction must be
naturally aligned: its first address is a multiple of its size, which
is the \term{alignment}[アライメント] requirement.\margin{CPG §5.3.2
Device Memory Accesses: \enquote{device memory is accessed via 32-, 64-,
or 128-byte memory transactions.  These memory transactions must be
naturally aligned} \cite{nvidia2024cudaguide}.}  Memory is cached in
units of a \term{cache line}[キャッシュライン], \SI{128}{\byte} on
NVIDIA GPUs, made of four \SI{32}{\byte} sectors, and a warp's accesses
cost as many sectors as they touch.  Finally the \term*{stride}
of an access, the distance between the addresses of consecutive threads,
decides how many sectors the warp's 32 addresses span.

\begin{example}
  \label{ex:l04-sectors}
  A warp of 32 threads loads one 4-byte word per thread,
  \SI{128}{\byte} of useful data, in the four patterns of
  \cref{fig:l04-sectors} and \cref{tab:l04-sectors}.  The efficiency is
  the useful data divided by the bytes transferred.
\end{example}

\begin{table}[htbp]
  \centering
  \caption{Sectors transferred for a warp loading 32 words of
    \SI{4}{\byte} (after BPG §10.2.1.1--10.2.1.4).}
  \label{tab:l04-sectors}
  \begin{tabular}{>{\raggedright\arraybackslash}p{52mm}ccc}
    \toprule
    access pattern & sectors & bytes moved & efficiency \\
    \midrule
    (a) consecutive, \SI{128}{\byte}-aligned & 4 & 128 & \SI{100}{\percent} \\
    (b) consecutive, offset by one word & 5 & 160 & \SI{80}{\percent} \\
    (c) stride 2 words & 8 & 256 & \SI{50}{\percent} \\
    (d) stride 8 words or more & 32 & 1024 & \SI{12.5}{\percent} \\
    \bottomrule
  \end{tabular}
\end{table}

\begin{figure}[htbp]
  \centering
  % Sectors touched by a 32-thread warp loading 4-byte words (Lesson 4).
% Each cell is one 4-byte word; thick rules mark 32-byte sectors.
% (a) consecutive and 128-byte aligned: 4 sectors; (b) the same shifted
% by one word: 5 sectors; (c) stride 2 words: 8 sectors; (d) stride 8
% words: one sector per thread.  Original figure (after BPG
% section 10.2.1.1--10.2.1.4).
% Usage: % Sectors touched by a 32-thread warp loading 4-byte words (Lesson 4).
% Each cell is one 4-byte word; thick rules mark 32-byte sectors.
% (a) consecutive and 128-byte aligned: 4 sectors; (b) the same shifted
% by one word: 5 sectors; (c) stride 2 words: 8 sectors; (d) stride 8
% words: one sector per thread.  Original figure (after BPG
% section 10.2.1.1--10.2.1.4).
% Usage: % Sectors touched by a 32-thread warp loading 4-byte words (Lesson 4).
% Each cell is one 4-byte word; thick rules mark 32-byte sectors.
% (a) consecutive and 128-byte aligned: 4 sectors; (b) the same shifted
% by one word: 5 sectors; (c) stride 2 words: 8 sectors; (d) stride 8
% words: one sector per thread.  Original figure (after BPG
% section 10.2.1.1--10.2.1.4).
% Usage: \input{figures/l04-sectors}
% Width: body (120 mm); the drawing is about 118 mm wide.
\begin{tikzpicture}[x=1.45mm, y=1mm, font=\scriptsize\sffamily]
  % \rowbase{y}{label}{sectors}: draws 8 cells per sector
  \newcommand{\rowbase}[3]{%
    \node[anchor=east, inner sep=1pt, align=right] at (-0.3,#1-2) {#2};
    \pgfmathtruncatemacro{\ncell}{8*#3-1}
    \foreach \c in {0,...,\ncell} {
      \draw[ln-rule] (\c,#1) rectangle ++(1,-4);
    }
    \foreach \s in {0,...,#3} {
      \draw[thick] (8*\s,#1+0.6) -- (8*\s,#1-4.6);
    }}
  % (a) aligned, consecutive
  \rowbase{30}{(a)}{8}
  \foreach \c in {0,...,31} { \draw[fill=ln-tip!40, ln-rule] (\c,30) rectangle ++(1,-4); }
  \node[anchor=west, text=ln-muted] at (64.5,28) {\iflangja{4セクタ，100\,\%}{4 sectors, \SI{100}{\percent}}};
  % (b) shifted by one word
  \rowbase{22}{(b)}{8}
  \foreach \c in {1,...,32} { \draw[fill=ln-tip!40, ln-rule] (\c,22) rectangle ++(1,-4); }
  \node[anchor=west, text=ln-muted] at (64.5,20) {\iflangja{5セクタ，80\,\%}{5 sectors, \SI{80}{\percent}}};
  % (c) stride 2
  \rowbase{14}{(c)}{8}
  \foreach \c in {0,2,...,62} { \draw[fill=ln-tip!40, ln-rule] (\c,14) rectangle ++(1,-4); }
  \node[anchor=west, text=ln-muted] at (64.5,12) {\iflangja{8セクタ，50\,\%}{8 sectors, \SI{50}{\percent}}};
  % (d) stride 8 (one word per sector)
  \rowbase{6}{(d)}{8}
  \foreach \c in {0,8,...,56} { \draw[fill=ln-tip!40, ln-rule] (\c,6) rectangle ++(1,-4); }
  \node[anchor=west, text=ln-muted, align=left] at (64.5,4) {\iflangja{$\cdots$ 32セクタ，12.5\,\%}{$\cdots$ 32 sectors, \SI{12.5}{\percent}}};
  % scale
  \draw[<->, ln-muted] (0,-1.5) -- (8,-1.5) node[midway, below, inner sep=1pt] {\SI{32}{\byte}};
  \draw[<->, ln-muted] (0,-5.5) -- (32,-5.5) node[midway, below, inner sep=1pt] {\SI{128}{\byte}\iflangja{（キャッシュライン）}{ (cache line)}};
\end{tikzpicture}

% Width: body (120 mm); the drawing is about 118 mm wide.
\begin{tikzpicture}[x=1.45mm, y=1mm, font=\scriptsize\sffamily]
  % \rowbase{y}{label}{sectors}: draws 8 cells per sector
  \newcommand{\rowbase}[3]{%
    \node[anchor=east, inner sep=1pt, align=right] at (-0.3,#1-2) {#2};
    \pgfmathtruncatemacro{\ncell}{8*#3-1}
    \foreach \c in {0,...,\ncell} {
      \draw[ln-rule] (\c,#1) rectangle ++(1,-4);
    }
    \foreach \s in {0,...,#3} {
      \draw[thick] (8*\s,#1+0.6) -- (8*\s,#1-4.6);
    }}
  % (a) aligned, consecutive
  \rowbase{30}{(a)}{8}
  \foreach \c in {0,...,31} { \draw[fill=ln-tip!40, ln-rule] (\c,30) rectangle ++(1,-4); }
  \node[anchor=west, text=ln-muted] at (64.5,28) {\iflangja{4セクタ，100\,\%}{4 sectors, \SI{100}{\percent}}};
  % (b) shifted by one word
  \rowbase{22}{(b)}{8}
  \foreach \c in {1,...,32} { \draw[fill=ln-tip!40, ln-rule] (\c,22) rectangle ++(1,-4); }
  \node[anchor=west, text=ln-muted] at (64.5,20) {\iflangja{5セクタ，80\,\%}{5 sectors, \SI{80}{\percent}}};
  % (c) stride 2
  \rowbase{14}{(c)}{8}
  \foreach \c in {0,2,...,62} { \draw[fill=ln-tip!40, ln-rule] (\c,14) rectangle ++(1,-4); }
  \node[anchor=west, text=ln-muted] at (64.5,12) {\iflangja{8セクタ，50\,\%}{8 sectors, \SI{50}{\percent}}};
  % (d) stride 8 (one word per sector)
  \rowbase{6}{(d)}{8}
  \foreach \c in {0,8,...,56} { \draw[fill=ln-tip!40, ln-rule] (\c,6) rectangle ++(1,-4); }
  \node[anchor=west, text=ln-muted, align=left] at (64.5,4) {\iflangja{$\cdots$ 32セクタ，12.5\,\%}{$\cdots$ 32 sectors, \SI{12.5}{\percent}}};
  % scale
  \draw[<->, ln-muted] (0,-1.5) -- (8,-1.5) node[midway, below, inner sep=1pt] {\SI{32}{\byte}};
  \draw[<->, ln-muted] (0,-5.5) -- (32,-5.5) node[midway, below, inner sep=1pt] {\SI{128}{\byte}\iflangja{（キャッシュライン）}{ (cache line)}};
\end{tikzpicture}

% Width: body (120 mm); the drawing is about 118 mm wide.
\begin{tikzpicture}[x=1.45mm, y=1mm, font=\scriptsize\sffamily]
  % \rowbase{y}{label}{sectors}: draws 8 cells per sector
  \newcommand{\rowbase}[3]{%
    \node[anchor=east, inner sep=1pt, align=right] at (-0.3,#1-2) {#2};
    \pgfmathtruncatemacro{\ncell}{8*#3-1}
    \foreach \c in {0,...,\ncell} {
      \draw[ln-rule] (\c,#1) rectangle ++(1,-4);
    }
    \foreach \s in {0,...,#3} {
      \draw[thick] (8*\s,#1+0.6) -- (8*\s,#1-4.6);
    }}
  % (a) aligned, consecutive
  \rowbase{30}{(a)}{8}
  \foreach \c in {0,...,31} { \draw[fill=ln-tip!40, ln-rule] (\c,30) rectangle ++(1,-4); }
  \node[anchor=west, text=ln-muted] at (64.5,28) {\iflangja{4セクタ，100\,\%}{4 sectors, \SI{100}{\percent}}};
  % (b) shifted by one word
  \rowbase{22}{(b)}{8}
  \foreach \c in {1,...,32} { \draw[fill=ln-tip!40, ln-rule] (\c,22) rectangle ++(1,-4); }
  \node[anchor=west, text=ln-muted] at (64.5,20) {\iflangja{5セクタ，80\,\%}{5 sectors, \SI{80}{\percent}}};
  % (c) stride 2
  \rowbase{14}{(c)}{8}
  \foreach \c in {0,2,...,62} { \draw[fill=ln-tip!40, ln-rule] (\c,14) rectangle ++(1,-4); }
  \node[anchor=west, text=ln-muted] at (64.5,12) {\iflangja{8セクタ，50\,\%}{8 sectors, \SI{50}{\percent}}};
  % (d) stride 8 (one word per sector)
  \rowbase{6}{(d)}{8}
  \foreach \c in {0,8,...,56} { \draw[fill=ln-tip!40, ln-rule] (\c,6) rectangle ++(1,-4); }
  \node[anchor=west, text=ln-muted, align=left] at (64.5,4) {\iflangja{$\cdots$ 32セクタ，12.5\,\%}{$\cdots$ 32 sectors, \SI{12.5}{\percent}}};
  % scale
  \draw[<->, ln-muted] (0,-1.5) -- (8,-1.5) node[midway, below, inner sep=1pt] {\SI{32}{\byte}};
  \draw[<->, ln-muted] (0,-5.5) -- (32,-5.5) node[midway, below, inner sep=1pt] {\SI{128}{\byte}\iflangja{（キャッシュライン）}{ (cache line)}};
\end{tikzpicture}

  \caption{Words touched by a 32-thread warp (shaded) over
    \SI{32}{\byte} sectors: (a) aligned and consecutive, (b) shifted by
    one word, (c) stride 2, (d) stride 8 (one word per sector).}
  \label{fig:l04-sectors}
\end{figure}

Pattern (b) shows why alignment matters even for consecutive accesses:
a shift of one word spills into a fifth sector.  Pattern (d) is the
column-wise access to a row-major matrix, and the uncoalesced load of
\cref{fig:l04-coalescing} is its extreme form; a kernel written so that
consecutive threads touch consecutive words, pattern (a), is the single
most important memory optimisation.\margin{K12 §1.4.3: GPUs before
compute capability 1.2 imposed strict coalescing rules (consecutive
threads, consecutive aligned words); later generations relaxed them and
produce \enquote{one to four memory transactions} for a coalesced warp
access \cite{kim2012gpgpu}.}  In CUDA the pattern is determined by how
the thread index enters the address:
\begin{lstlisting}[language=C++, numbers=none]
// coalesced: consecutive threads, consecutive words
out[i] = in[blockIdx.x * blockDim.x + threadIdx.x];
// uncoalesced: stride of one row (n words) per thread
out[i] = in[threadIdx.x * n + blockIdx.x];
\end{lstlisting}
\Needspace{18\baselineskip}
When the algorithm needs the strided view, the remedy is the shared
memory of \cref{sec:l04-bank-conflicts}: load a tile coalesced, then
read it from shared memory in the order the computation wants, which is
the tiling of Lesson~3 and the subject of Lesson~5.\tip{Check a load by
asking what addresses threads 0 and 1 of a warp touch: if they differ
by the element size and thread 0's address is aligned, the access is
coalesced.}

\begin{keypoint}[Lesson summary]
  An SM hides memory latency with fine-grained multithreading: one PC
  per warp and a register set per thread stay on chip, so a warp switch
  is free and the number of warps is bounded by storage.  The register
  file and shared memory are banked to supply a warp's operands per
  cycle; accesses that need more values from one bank than it has ports
  are bank conflicts, absorbed in the register file by operand
  collectors and a scoreboard that issues whichever warp is ready, and
  avoided in shared memory (32 banks of \SI{4}{\byte}) by odd strides
  and padding.  A warp instruction is fetched and decoded once and
  executed per lane under an active mask, with immediates and
  \texttt{\%ctaid.x} broadcast and \texttt{\%tid.x} per lane.  On the
  memory side one warp load can issue 32 requests; the hardware
  coalesces those that fall into the same aligned \SI{32}{\byte}
  sectors, so consecutive, aligned accesses move 4 sectors for
  \SI{128}{\byte} of data while strided ones move up to 32, and HBM's
  stacked DRAM on an interposer exists to serve the remaining demand.
\end{keypoint}

\end{document}
